\input{preamble}

% OK, start here.
%
\begin{document}

\title{Descente et espaces algébriques}

\maketitle

\phantomsection
\label{section-phantom}

\tableofcontents

\section{Introduction}
\label{section-introduction}

\noindent
Dans le chapitre consacré aux topologies sur les espaces algébriques (voir
Topologies sur les espaces, section \ref{spaces-topologies-section-introduction})
nous avons introduit les recouvrements \'etales, fppf, lisses, syntomiques et fpqc des
espaces algébriques.
Dans ce chapitre, nous étudions les structures sur les espaces algébriques
dont la descente est possible relativement à de tels recouvrements.
Voir par exemple \cite{Gr-I}, \cite{Gr-II}, \cite{Gr-III},
\cite{Gr-IV}, \cite{Gr-V} et \cite{Gr-VI}.



\section{Conventions}
\label{section-conventions}

\noindent
L'hypothèse permanente est que tous les schémas appartiennent à
un grand site fppf $\Sch_{fppf}$. De plus, tous les anneaux $A$ considérés
sont tels que $\Spec(A)$ est isomorphe à un
objet de ce grand site.

\medskip\noindent
Soit $S$ un schéma et soit $X$ un espace algébrique sur $S$.
Dans ce chapitre et dans le suivant, nous écrirons $X \times_S X$
pour le produit de $X$ avec lui-même (dans la catégorie des espaces algébriques
sur $S$), au lieu de $X \times X$.






\section{Données de descente pour les faisceaux quasi-cohérents}
\label{section-equivalence}

\noindent
Cette section est l'analogue de
Descente, section \ref{descent-section-equivalence}
pour les espaces algébriques.
Il est naturel de lire d'abord cette dernière.

\begin{definition}
\label{definition-descent-datum-quasi-coherent}
Soit $S$ un schéma. Soit $\{f_i : X_i \to X\}_{i \in I}$ une famille
de morphismes d'espaces algébriques sur $S$ à cible fixée $X$.
\begin{enumerate}
\item Une {\it donnée de descente $(\mathcal{F}_i, \varphi_{ij})$
pour les faisceaux quasi-cohérents} relativement à la famille donnée
consiste en la donnée d'un faisceau quasi-cohérent $\mathcal{F}_i$ sur $X_i$ pour
chaque $i \in I$, et d'un isomorphisme de
$\mathcal{O}_{X_i \times_X X_j}$-modules quasi-cohérents
$\varphi_{ij} : \text{pr}_0^*\mathcal{F}_i \to \text{pr}_1^*\mathcal{F}_j$
pour chaque paire $(i, j) \in I^2$
tel que, pour tout triplet d'indices $(i, j, k) \in I^3$, le
diagramme
$$
\xymatrix{
\text{pr}_0^*\mathcal{F}_i \ar[rd]_{\text{pr}_{01}^*\varphi_{ij}}
\ar[rr]_{\text{pr}_{02}^*\varphi_{ik}} & &
\text{pr}_2^*\mathcal{F}_k \\
& \text{pr}_1^*\mathcal{F}_j \ar[ru]_{\text{pr}_{12}^*\varphi_{jk}} &
}
$$
de $\mathcal{O}_{X_i \times_X X_j \times_X X_k}$-modules
est commutatif. C'est la {\it condition de cocycle}.
\item Un {\it morphisme $\psi : (\mathcal{F}_i, \varphi_{ij}) \to
(\mathcal{F}'_i, \varphi'_{ij})$ de données de descente} est donné
par une famille $\psi = (\psi_i)_{i\in I}$ de morphismes de
$\mathcal{O}_{X_i}$-modules $\psi_i : \mathcal{F}_i \to \mathcal{F}'_i$
telle que tous les diagrammes
$$
\xymatrix{
\text{pr}_0^*\mathcal{F}_i \ar[r]_{\varphi_{ij}} \ar[d]_{\text{pr}_0^*\psi_i}
& \text{pr}_1^*\mathcal{F}_j \ar[d]^{\text{pr}_1^*\psi_j} \\
\text{pr}_0^*\mathcal{F}'_i \ar[r]^{\varphi'_{ij}} &
\text{pr}_1^*\mathcal{F}'_j \\
}
$$
soient commutatifs.
\end{enumerate}
\end{definition}

\begin{lemma}
\label{lemma-map-families}
Soit $S$ un schéma.
Soient $\mathcal{U} = \{U_i \to U\}_{i \in I}$ et
$\mathcal{V} = \{V_j \to V\}_{j \in J}$
des familles de morphismes d'espaces algébriques sur $S$ à cibles fixées.
Soit $(g, \alpha : I \to J, (g_i)) : \mathcal{U} \to \mathcal{V}$
un morphisme de familles de morphismes à cible fixée, voir
Sites, définition \ref{sites-definition-morphism-coverings}.
Soit $(\mathcal{F}_j, \varphi_{jj'})$ une donnée de descente
pour les faisceaux quasi-cohérents relativement à la
famille $\{V_j \to V\}_{j \in J}$. Alors
\begin{enumerate}
\item Le système
$$
\left(g_i^*\mathcal{F}_{\alpha(i)},
(g_i \times g_{i'})^*\varphi_{\alpha(i)\alpha(i')}\right)
$$
est une donnée de descente relativement à la famille $\{U_i \to U\}_{i \in I}$.
\item Cette construction est fonctorielle en la donnée de descente
$(\mathcal{F}_j, \varphi_{jj'})$.
\item Étant donné un second morphisme
$(g', \alpha' : I \to J, (g'_i))$
de familles de morphismes à cible fixée
tel que $g = g'$, il existe un isomorphisme fonctoriel de données de descente
$$
(g_i^*\mathcal{F}_{\alpha(i)},
(g_i \times g_{i'})^*\varphi_{\alpha(i)\alpha(i')})
\cong
((g'_i)^*\mathcal{F}_{\alpha'(i)},
(g'_i \times g'_{i'})^*\varphi_{\alpha'(i)\alpha'(i')}).
$$
\end{enumerate}
\end{lemma}

\begin{proof}
Démonstration omise. Indication : les morphismes
$g_i^*\mathcal{F}_{\alpha(i)} \to (g'_i)^*\mathcal{F}_{\alpha'(i)}$
qui donnent l'isomorphisme de données de descente en (3)
sont les images réciproques des morphismes $\varphi_{\alpha(i)\alpha'(i)}$ par les
morphismes $(g_i, g'_i) : U_i \to V_{\alpha(i)} \times_V V_{\alpha'(i)}$.
\end{proof}

\noindent
Soit $g : U \to V$ un morphisme d'espaces algébriques.
Le lemme ci-dessus montre qu'il existe un foncteur image réciproque bien défini
entre les catégories de données de descente relatives à des familles de
morphismes ayant respectivement $V$ et $U$ pour buts, pourvu qu'il existe entre ces
familles un morphisme qui ``vit au-dessus de $g$''.

\begin{definition}
\label{definition-descent-datum-effective-quasi-coherent}
Soit $S$ un schéma.
Soit $\{U_i \to U\}_{i \in I}$ une famille de morphismes d'espaces algébriques
sur $S$ à cible fixée.
\begin{enumerate}
\item Soit $\mathcal{F}$ un $\mathcal{O}_U$-module quasi-cohérent.
L'unique donnée de descente portée par $\mathcal{F}$ relativement au recouvrement
$\{U \to U\}$ est appelée la {\it donnée de descente triviale}.
\item L'image réciproque de la donnée de descente triviale relativement à
la famille $\{U_i \to U\}$ est appelée la {\it donnée de descente canonique}.
Notation : $(\mathcal{F}|_{U_i}, can)$.
\item Une donnée de descente $(\mathcal{F}_i, \varphi_{ij})$
pour les faisceaux quasi-cohérents relativement à la famille donnée
est dite {\it effective} s'il existe un faisceau quasi-cohérent
$\mathcal{F}$ sur $U$ tel que $(\mathcal{F}_i, \varphi_{ij})$
soit isomorphe à $(\mathcal{F}|_{U_i}, can)$.
\end{enumerate}
\end{definition}

\begin{lemma}
\label{lemma-zariski-descent-effective}
Soit $S$ un schéma. Soit $U$ un espace algébrique sur $S$.
Soit $\{U_i \to U\}$ un recouvrement de Zariski de $U$, voir
Topologies sur les espaces,
définition \ref{spaces-topologies-definition-zariski-covering}.
Toute donnée de descente pour les faisceaux quasi-cohérents
relativement à la famille $\mathcal{U} = \{U_i \to U\}$ est
effective. De plus, le foncteur de la catégorie des
$\mathcal{O}_U$-modules quasi-cohérents vers la catégorie
des données de descente relativement à $\{U_i \to U\}$ est pleinement fidèle.
\end{lemma}

\begin{proof}
Démonstration omise.
\end{proof}












\section{Descente fpqc des faisceaux quasi-cohérents}
\label{section-fpqc-descent-quasi-coherent}

\noindent
La principale application de la descente plate des modules est
l'énoncé de descente correspondant pour les faisceaux quasi-cohérents
relativement aux recouvrements fpqc.

\begin{proposition}
\label{proposition-fpqc-descent-quasi-coherent}
Soit $S$ un schéma.
Soit $\{X_i \to X\}$ un recouvrement fpqc d'espaces algébriques sur $S$, voir
Topologies sur les espaces,
définition \ref{spaces-topologies-definition-fpqc-covering}.
Toute donnée de descente pour les faisceaux quasi-cohérents
relativement à $\{X_i \to X\}$ est effective.
De plus, le foncteur de la catégorie des
$\mathcal{O}_X$-modules quasi-cohérents vers la catégorie
des données de descente relativement à $\{X_i \to X\}$ est pleinement fidèle.
\end{proposition}

\begin{proof}
C'est, à peu de chose près, une conséquence formelle du
résultat correspondant pour les schémas, voir
Descente, proposition \ref{descent-proposition-fpqc-descent-quasi-coherent}.
Voici une stratégie de démonstration :
\begin{enumerate}
\item Le fait que $\{X_i \to X\}$ soit un raffinement du recouvrement trivial
$\{X \to X\}$ donne, via
le lemme \ref{lemma-map-families},
un foncteur $\QCoh(\mathcal{O}_X) \to DD(\{X_i \to X\})$ de la
catégorie des $\mathcal{O}_X$-modules quasi-cohérents vers la catégorie des
données de descente pour la famille donnée.
\item Afin de démontrer la proposition, nous allons construire un
foncteur quasi-inverse
$back : DD(\{X_i \to X\}) \to \QCoh(\mathcal{O}_X)$.
\item En appliquant à nouveau
le lemme \ref{lemma-map-families}
nous voyons qu'il existe un foncteur
$DD(\{X_i \to X\}) \to DD(\{T_j \to X\})$
si $\{T_j \to X\}$ est un raffinement de la famille donnée.
Ainsi, pour construire le foncteur $back$, nous pouvons supposer que
chaque $X_i$ est un schéma, voir
Topologies sur les espaces,
lemme \ref{spaces-topologies-lemma-refine-fpqc-schemes}.
Cela nous ramène au cas où tous les $X_i$ sont des schémas.
\item Un faisceau quasi-cohérent sur $X$ est, par définition, un
$\mathcal{O}_X$-module quasi-cohérent sur $X_\etale$. Pour tout
$U \in \Ob(X_\etale)$, on obtient alors un recouvrement fppf
$\{U_i \times_X X_i \to U\}$ par des schémas et un morphisme
$g : \{U_i \times_X X_i \to U\} \to \{X_i \to X\}$ de recouvrements
au-dessus de $U \to X$. Étant donnée une donnée de descente
$\xi = (\mathcal{F}_i, \varphi_{ij})$, on obtient un
$\mathcal{O}_U$-module quasi-cohérent $\mathcal{F}_{\xi, U}$ correspondant
à l'image réciproque $g^*\xi$ fournie par le
lemme \ref{lemma-map-families}
sur le recouvrement de $U$, puis on applique l'effectivité de la descente
pour les recouvrements fppf de schémas, voir Descente, proposition \ref{descent-proposition-fpqc-descent-quasi-coherent}.
\item Vérifier que $\xi \mapsto \mathcal{F}_{\xi, U}$ est fonctorielle en $\xi$.
Vérification omise.
\item Vérifier que $\xi \mapsto \mathcal{F}_{\xi, U}$ est compatible
avec les morphismes $U \to U'$ du site $X_\etale$, de sorte que
le système de faisceaux $\mathcal{F}_{\xi, U}$ corresponde à un faisceau quasi-cohérent
$\mathcal{F}_\xi$ sur $X_\etale$, voir
Propriétés des espaces,
lemme \ref{spaces-properties-lemma-characterize-quasi-coherent-small-etale}.
Détails omis.
\item Vérifier que $back : \xi \mapsto \mathcal{F}_\xi$ est quasi-inverse
du foncteur construit en (1). Vérification omise.
\end{enumerate}
La démonstration est ainsi achevée.
\end{proof}









\section{Modules quasi-cohérents et schémas affines}
\label{section-alternative-quasi-coherent}

\noindent
Soit $S$ un schéma. Soit $X$ un espace algébrique sur $S$.
Rappelons que $X_{affine, \etale}$ est la sous-catégorie pleine de $X_\etale$
dont les objets sont les schémas affines; on la munit d'une structure de site en prenant
pour recouvrements les recouvrements \'etales standards. Voir Propriétés des espaces, définition
\ref{spaces-properties-definition-affine-etale-site}.
D'après Propriétés des espaces, lemme \ref{spaces-properties-lemma-alternative},
on a une équivalence de topos
$g : \Sh(X_{affine, \etale}) \to \Sh(X_\etale)$
dont le foncteur image réciproque est donné par la restriction.
Rappelons que $\mathcal{O}_X$ désigne le faisceau structural sur
$X_\etale$. Nous obtenons alors une équivalence
\begin{equation}
\label{equation-alternative-small-ringed}
(\Sh(X_{affine, \etale}), \mathcal{O}_X|_{X_{affine, \etale}})
\longrightarrow
(\Sh(X_\etale), \mathcal{O}_X)
\end{equation}
de topos annelés. Nous écrirons souvent $\mathcal{O}_X$
au lieu de $\mathcal{O}_X|_{X_{affine, \etale}}$.
Cela permet également de comparer les modules quasi-cohérents.

\begin{lemma}
\label{lemma-quasi-coherent-alternative-small}
Soit $S$ un schéma. Soit $X$ un espace algébrique sur $S$.
Soit $\mathcal{F}$ un préfaisceau de $\mathcal{O}_X$-modules
sur $X_{affine, \etale}$. Les conditions suivantes sont équivalentes :
\begin{enumerate}
\item pour tout morphisme $U \to U'$ de $X_{affine, \etale}$ la flèche
$\mathcal{F}(U') \otimes_{\mathcal{O}_X(U')} \mathcal{O}_X(U)
\to \mathcal{F}(U)$ est un isomorphisme,
\item $\mathcal{F}$ est un module quasi-cohérent sur le site annelé
$(X_{affine, \etale}, \mathcal{O}_X)$ au sens de
Modules sur les sites, définition \ref{sites-modules-definition-site-local},
\item $\mathcal{F}$ correspond à un module quasi-cohérent sur $X$
via l'équivalence (\ref{equation-alternative-small-ringed}),
\end{enumerate}
\end{lemma}

\begin{proof}
Supposons que (1) soit satisfaite. Pour montrer que $\mathcal{F}$ est un faisceau, soit
$\mathcal{U} = \{U_i \to U\}_{i = 1, \ldots, n}$ un recouvrement
de $X_{affine, \etale}$. La condition de faisceau pour $\mathcal{F}$
et $\mathcal{U}$, par notre hypothèse sur $\mathcal{F}$,
se réduit à montrer que
$$
0 \to \mathcal{F}(U) \to
\prod \mathcal{F}(U) \otimes_{\mathcal{O}_X(U)} \mathcal{O}_X(U_i) \to
\prod \mathcal{F}(U) \otimes_{\mathcal{O}_X(U)} \mathcal{O}_X(U_i \times_U U_j)
$$
Cette suite est exacte, car $\mathcal{O}_X(U) \to \prod \mathcal{O}_X(U_i)$
est fidèlement plat
(d'après Descente, lemme \ref{descent-lemma-standard-covering-Cech}, et parce que
les recouvrements de $X_{affine, \etale}$ sont les recouvrements \'etales
standards); on peut alors appliquer Descente, lemme \ref{descent-lemma-ff-exact}.
Ensuite, nous montrons que $\mathcal{F}$ est quasi-cohérent sur $X_{affine, \etale}$.
En effet, pour $U$ dans $X_{affine, \etale}$, posons $R = \mathcal{O}_X(U)$
et choisissons une présentation
$$
\bigoplus\nolimits_{k \in K} R
\longrightarrow
\bigoplus\nolimits_{l \in L} R
\longrightarrow
\mathcal{F}(U)
\longrightarrow 0
$$
par des $R$-modules libres. D'après la propriété (1) et l'exactitude à droite du produit tensoriel,
nous voyons que, pour tout morphisme $U' \to U$ dans $X_{affine, \etale}$,
on obtient une présentation
$$
\bigoplus\nolimits_{k \in K} \mathcal{O}_X(U')
\longrightarrow
\bigoplus\nolimits_{l \in L} \mathcal{O}_X(U')
\longrightarrow
\mathcal{F}(U')
\longrightarrow 0
$$
Autrement dit, la restriction de $\mathcal{F}$
à la catégorie localisée $X_{affine, etale}/U$ possède une présentation
$$
\bigoplus\nolimits_{k \in K} \mathcal{O}_X|_{X_{affine, \etale}/U}
\longrightarrow
\bigoplus\nolimits_{l \in L} \mathcal{O}_X|_{X_{affine, \etale}/U}
\longrightarrow
\mathcal{F}|_{X_{affine, \etale}/U}
\longrightarrow 0
$$
ce qui est la présentation requise pour montrer que $\mathcal{F}$ est quasi-cohérent.
Toutes nos excuses pour cette horrible notation; cela achève la démonstration
que (1) implique (2).

\medskip\noindent
Puisque la notion de module quasi-cohérent est intrinsèque
(Modules sur les sites, lemme \ref{sites-modules-lemma-special-locally-free})
nous voyons que l'équivalence (\ref{equation-alternative-small-ringed})
induit une équivalence entre les catégories de modules quasi-cohérents.
Ainsi, (2) et (3) sont équivalentes.

\medskip\noindent
Supposons (3) et démontrons (1). Soit
$\mathcal{G}$ un module quasi-cohérent sur $X$ correspondant à $\mathcal{F}$.
Soit $h : U \to U' \to X$ un morphisme de $X_{affine, \etale}$.
Notons $f : U \to X$ et $f' : U' \to X$ les morphismes structuraux,
de sorte que $f = f' \circ h$. Nous avons
$\mathcal{F}(U') = \Gamma(U', (f')^*\mathcal{G})$ et
$\mathcal{F}(U) = \Gamma(U, f^*\mathcal{G}) = \Gamma(U, h^*(f')^*\mathcal{G})$.
Ainsi, (1) résulte de Schémas, lemme \ref{schemes-lemma-widetilde-pullback}.
\end{proof}
% preserved blank authority line 391
% preserved blank authority line 392
% preserved blank authority line 393
% preserved blank authority line 394
% preserved blank authority line 395
\section{Descente des propriétés de finitude des modules}
\label{section-descent-finiteness}

\noindent
Cette section est, pour les espaces algébriques, l'analogue de
Descente, section \ref{descent-section-descent-finiteness}.
Nous voulons montrer que certaines conditions de finitude portant sur un module
quasi-cohérent peuvent se vérifier sur les membres d'un
recouvrement. Nous utiliserons à plusieurs reprises le schéma de démonstration suivant.
Supposons que $X$ soit un espace algébrique et que $\{X_i \to X\}$
soit un recouvrement fppf (resp.\ fpqc). Soit $U \to X$ un morphisme \'etale
surjectif tel que $U$ soit un schéma. Il existe alors un
recouvrement fppf (resp.\ fpqc) $\{Y_j \to X\}$ tel que
\begin{enumerate}
\item $\{Y_j \to X\}$ est un raffinement de $\{X_i \to X\}$,
\item chaque $Y_j$ est un schéma, et
\item chaque morphisme $Y_j \to X$ se factorise par $U$, et
\item $\{Y_j \to U\}$ est un recouvrement fppf (resp.\ fpqc) de $U$.
\end{enumerate}
En effet, raffinons d'abord $\{X_i \to X\}$ par un recouvrement fppf (resp.\ fpqc)
tel que chaque $X_i$ soit un schéma, voir
Topologies sur les espaces, lemme \ref{spaces-topologies-lemma-refine-fppf-schemes},
resp.\ le lemme \ref{spaces-topologies-lemma-refine-fpqc-schemes}.
Posons alors $Y_i = U \times_X X_i$. Un module quasi-cohérent
sur $\mathcal{O}_X$, noté $\mathcal{F}$, est de type fini, de
présentation finie, etc., si et seulement si le module quasi-cohérent
sur $\mathcal{O}_U$, noté $\mathcal{F}|_U$, est de type fini, de
présentation finie, etc. L'existence du
raffinement $\{Y_j \to X\}$ permet donc de ramener la démonstration des
lemmes suivants au cas des schémas. Nous l'indiquerons en disant
que « {\it le résultat se déduit du cas des schémas par
localisation \'etale} ».

\begin{lemma}
\label{lemma-finite-type-descends}
Soit $X$ un espace algébrique sur un schéma $S$.
Soit $\mathcal{F}$ un $\mathcal{O}_X$-module quasi-cohérent.
Soit $\{f_i : X_i \to X\}_{i \in I}$ un recouvrement fpqc tel que
chaque $f_i^*\mathcal{F}$ soit un $\mathcal{O}_{X_i}$-module de type fini.
Alors $\mathcal{F}$ est un $\mathcal{O}_X$-module de type fini.
\end{lemma}

\begin{proof}
Ce résultat se déduit du cas des schémas, voir
Descente, lemme \ref{descent-lemma-finite-type-descends},
par localisation \'etale.
\end{proof}

\begin{lemma}
\label{lemma-finite-presentation-descends}
Soit $X$ un espace algébrique sur un schéma $S$.
Soit $\mathcal{F}$ un $\mathcal{O}_X$-module quasi-cohérent.
Soit $\{f_i : X_i \to X\}_{i \in I}$ un recouvrement fpqc tel que
chaque $f_i^*\mathcal{F}$ soit un $\mathcal{O}_{X_i}$-module de présentation
finie. Alors $\mathcal{F}$ est un $\mathcal{O}_X$-module
de présentation finie.
\end{lemma}

\begin{proof}
Ce résultat se déduit du cas des schémas, voir
Descente, lemme \ref{descent-lemma-finite-presentation-descends},
par localisation \'etale.
\end{proof}

\begin{lemma}
\label{lemma-flat-descends}
Soit $X$ un espace algébrique sur un schéma $S$.
Soit $\mathcal{F}$ un $\mathcal{O}_X$-module quasi-cohérent.
Soit $\{f_i : X_i \to X\}_{i \in I}$ un recouvrement fpqc tel que
chaque $f_i^*\mathcal{F}$ soit un $\mathcal{O}_{X_i}$-module plat.
Alors $\mathcal{F}$ est un $\mathcal{O}_X$-module plat.
\end{lemma}

\begin{proof}
Ce résultat se déduit du cas des schémas, voir
Descente, lemme \ref{descent-lemma-flat-descends},
par localisation \'etale.
\end{proof}

\begin{lemma}
\label{lemma-finite-locally-free-descends}
Soit $X$ un espace algébrique sur un schéma $S$.
Soit $\mathcal{F}$ un $\mathcal{O}_X$-module quasi-cohérent.
Soit $\{f_i : X_i \to X\}_{i \in I}$ un recouvrement fpqc tel que
chaque $f_i^*\mathcal{F}$ soit un $\mathcal{O}_{X_i}$-module localement libre de type fini.
Alors $\mathcal{F}$ est un $\mathcal{O}_X$-module localement libre de type fini.
\end{lemma}

\begin{proof}
Ce résultat se déduit du cas des schémas, voir
Descente, lemme \ref{descent-lemma-finite-locally-free-descends},
par localisation \'etale.
\end{proof}

\noindent
La définition d'un faisceau quasi-cohérent localement projectif se trouve dans
Propriétés des espaces, section
\ref{spaces-properties-section-locally-projective}.
Il y est également démontré que cette notion est stable par image réciproque.

\begin{lemma}
\label{lemma-locally-projective-descends}
Soit $X$ un espace algébrique sur un schéma $S$.
Soit $\mathcal{F}$ un $\mathcal{O}_X$-module quasi-cohérent.
Soit $\{f_i : X_i \to X\}_{i \in I}$ un recouvrement fpqc tel que
chaque $f_i^*\mathcal{F}$ soit un $\mathcal{O}_{X_i}$-module localement projectif.
Alors $\mathcal{F}$ est un $\mathcal{O}_X$-module localement projectif.
\end{lemma}

\begin{proof}
Ce résultat se déduit du cas des schémas, voir
Descente, lemme \ref{descent-lemma-locally-projective-descends},
par localisation \'etale.
\end{proof}

\noindent
Ajoutons ici deux résultats liés aux précédents, mais
de nature légèrement différente.

\begin{lemma}
\label{lemma-finite-over-finite-module}
Soit $S$ un schéma.
Soit $f : X \to Y$ un morphisme d'espaces algébriques sur $S$.
Soit $\mathcal{F}$ un $\mathcal{O}_X$-module quasi-cohérent.
Supposons que $f$ soit fini.
Alors $\mathcal{F}$ est un $\mathcal{O}_X$-module de type fini
si et seulement si $f_*\mathcal{F}$ est un $\mathcal{O}_Y$-module de type
fini.
\end{lemma}

\begin{proof}
Puisque $f$ est fini, il est représentable. Choisissons un schéma $V$ et un
morphisme \'etale surjectif $V \to Y$. Alors $U = V \times_Y X$ est un schéma muni d'un
morphisme \'etale surjectif vers $X$ et d'un morphisme fini
$\psi : U \to V$ (le changement de base de $f$). Puisque
$\psi_*(\mathcal{F}|_U) = f_*\mathcal{F}|_V$,
l'assertion se déduit immédiatement de la version pour les schémas,
c'est-à-dire du
lemme \ref{descent-lemma-finite-over-finite-module} de Descente.
\end{proof}

\begin{lemma}
\label{lemma-finite-finitely-presented-module}
Soit $S$ un schéma.
Soit $f : X \to Y$ un morphisme d'espaces algébriques sur $S$.
Soit $\mathcal{F}$ un $\mathcal{O}_X$-module quasi-cohérent.
Supposons que $f$ soit fini et de présentation finie.
Alors $\mathcal{F}$ est un $\mathcal{O}_X$-module de présentation finie
si et seulement si $f_*\mathcal{F}$ est un $\mathcal{O}_Y$-module de présentation
finie.
\end{lemma}

\begin{proof}
Puisque $f$ est fini, il est représentable. Choisissons un schéma $V$ et un
morphisme \'etale surjectif $V \to Y$. Alors $U = V \times_Y X$ est un schéma muni d'un
morphisme \'etale surjectif vers $X$ et d'un morphisme fini
$\psi : U \to V$ (le changement de base de $f$). Puisque
$\psi_*(\mathcal{F}|_U) = f_*\mathcal{F}|_V$,
l'assertion se déduit immédiatement de la version pour les schémas,
c'est-à-dire du
lemme \ref{descent-lemma-finite-finitely-presented-module} de Descente.
\end{proof}






\section{Recouvrements fpqc}
\label{section-fpqc}

\noindent
Cette section est l'analogue de la
section \ref{descent-section-fpqc-universal-effective-epimorphisms} de Descente.
Nous ne savons pas actuellement si tous les résultats valables pour les
recouvrements fpqc de schémas le sont aussi pour les espaces algébriques.

\begin{lemma}
\label{lemma-open-fpqc-covering}
Soit $S$ un schéma.
Soit $\{f_i : T_i \to T\}_{i \in I}$ un recouvrement fpqc
d'espaces algébriques sur $S$.
Pour chaque $i$, soit $W_i \subset T_i$ un sous-espace ouvert,
et supposons que, pour tous $i, j \in I$,
$\text{pr}_0^{-1}(W_i) = \text{pr}_1^{-1}(W_j)$ en tant que sous-espaces ouverts
de $T_i \times_T T_j$. Alors il existe un unique sous-espace ouvert
$W \subset T$ tel que $W_i = f_i^{-1}(W)$ pour chaque $i$.
\end{lemma}

\begin{proof}
D'après le
lemme \ref{spaces-topologies-lemma-refine-fpqc-schemes} de Topologies sur les espaces,
nous pouvons supposer que chaque $T_i$ est un schéma.
Choisissons un schéma $U$ et un morphisme \'etale surjectif $U \to T$.
Alors $\{T_i \times_T U \to U\}$ est un recouvrement fpqc de $U$
et $T_i \times_T U$ est un schéma pour chaque $i$. La famille
des ouverts $W_i \times_T U$ provient donc d'un unique
sous-schéma ouvert $W' \subset U$, d'après le
lemme \ref{descent-lemma-open-fpqc-covering} de Descente.
Comme $U \to X$ est ouvert, nous pouvons définir $W \subset X$ comme l'ouvert de Zariski
image de $W'$; voir
Propriétés des espaces, section \ref{spaces-properties-section-points}.
Nous omettons de vérifier que cette construction convient, c'est-à-dire que
$W_i$ est l'image réciproque de $W$ pour chaque $i$.
\end{proof}

\begin{lemma}
\label{lemma-fpqc-universal-effective-epimorphisms}
Soit $S$ un schéma. Soit $\{T_i \to T\}$ un recouvrement fpqc d'espaces algébriques
sur $S$; voir Topologies sur les espaces, définition
\ref{spaces-topologies-definition-fpqc-covering}.
Alors, pour tout espace algébrique $B$ sur $S$, la suite
$$
\xymatrix{
\Mor_S(T, B) \ar[r] &
\prod\nolimits_i \Mor_S(T_i, B) \ar@<1ex>[r] \ar@<-1ex>[r] &
\prod\nolimits_{i, j} \Mor_S(T_i \times_T T_j, B)
}
$$
est un diagramme égalisateur.
Autrement dit, tout foncteur représentable sur la catégorie des
espaces algébriques sur $S$ vérifie la condition de faisceau pour les
recouvrements fpqc.
\end{lemma}

\begin{proof}
Nous savons que l'assertion est vraie si $\{T_i \to T\}$ est un recouvrement fpqc de
schémas, voir Propriétés des espaces, proposition
\ref{spaces-properties-proposition-sheaf-fpqc}.
C'est le point essentiel; nous invitons le lecteur à omettre la suite
de cette démonstration, qui est formelle. Choisissons un schéma $U$ et un
morphisme \'etale surjectif
$U \to T$. Choisissons un schéma $U_i$ et un morphisme $U_i \to T_i \times_T U$
\'etale surjectif. Alors $\{U_i \to U\}$ est un
recouvrement fpqc. Cela découle de
Topologies sur les espaces, lemmes \ref{spaces-topologies-lemma-fpqc} et
\ref{spaces-topologies-lemma-recognize-fpqc-covering}.
Ce qui précède donne le résultat pour $\{U_i \to U\}$.

\medskip\noindent
Concrètement, supposons que les $b_i : T_i \to B$
forment une famille de morphismes telle que
$b_i \circ \text{pr}_0 = b_j \circ \text{pr}_1$ comme morphismes
$T_i \times_T T_j \to B$. Définissons $a_i : U_i \to B$ comme le
composé obtenu en faisant suivre $U_i \to T_i$ de $b_i$. D'après ce qui précède,
il existe un unique morphisme $a : U \to B$ tel que
$a_i$ soit le composé de $a$ avec $U_i \to U$.
L'unicité assure que $a \circ \text{pr}_0 = a \circ \text{pr}_1$
comme morphismes $U \times_T U \to B$. Puisque $T = U/(U \times_T U)$
comme faisceau, $a$ provient d'un unique morphisme $b : T \to B$.
Une chasse au diagramme montre que $b$ est le morphisme recherché.
\end{proof}










\section{Descente des propriétés de finitude et de lissité des morphismes}
\sectionmark{Descente : finitude et lissité}
\label{section-descent-finiteness-morphisms}

\noindent
Un lemme du type suivant est parfois utile.

\begin{lemma}
\label{lemma-curiosity}
Soit $S$ un schéma. Soient $X \to Y \to Z$ des morphismes d'espaces algébriques.
Soit $P$ l'une des propriétés suivantes pour un morphisme d'espaces algébriques
sur $S$ :
être plat, être localement de type fini ou être localement de présentation finie.
Supposons que $X \to Z$ vérifie $P$ et que
$X \to Y$ soit une surjection de faisceaux sur $(\Sch/S)_{fppf}$.
Alors $Y \to Z$ vérifie $P$.
\end{lemma}

\begin{proof}
Choisissons un schéma $W$ et un morphisme \'etale surjectif $W \to Z$.
Choisissons un schéma $V$ et un morphisme \'etale surjectif $V \to W \times_Z Y$.
Choisissons un schéma $U$ et un morphisme \'etale surjectif $U \to V \times_Y X$.
Par hypothèse, il existe un recouvrement fppf $\{V_i \to V\}$ et des
relèvements $V_i \to X$ du morphisme $V_i \to Y$. Comme $U \to X$ est surjectif
\'etale, nous voyons que, sur les membres du recouvrement fppf
$\{V_i \times_X U \to V\}$, il existe des relèvements dans $U$. Par conséquent, $U \to V$ induit
une surjection de faisceaux sur $(\Sch/S)_{fppf}$.
D'après la définition de la propriété $P$ pour un
morphisme d'espaces algébriques (voir
Morphismes d'espaces,
définition \ref{spaces-morphisms-definition-flat},
définition \ref{spaces-morphisms-definition-locally-finite-type}, et
définition \ref{spaces-morphisms-definition-locally-finite-presentation})
nous savons que $U \to W$ vérifie $P$, et il reste à montrer que $V \to W$ vérifie $P$.
La question est ainsi ramenée au cas des morphismes de schémas
traité dans le
lemme \ref{descent-lemma-curiosity} de Descente.
\end{proof}

\noindent
Un cas plus usuel du lemme précédent est le suivant.
(La version « plate » découle de
Morphismes d'espaces, lemme \ref{spaces-morphisms-lemma-flat-permanence}.)

\begin{lemma}
\label{lemma-flat-finitely-presented-permanence}
Soit $S$ un schéma. Considérons le diagramme suivant
$$
\xymatrix{
X \ar[rr]_f \ar[rd]_p & &
Y \ar[dl]^q \\
& B
}
$$
de morphismes d'espaces algébriques sur $S$, supposé commutatif.
Supposons que $f$ soit surjectif, plat et localement de présentation finie
et que $p$ soit localement de présentation finie (resp.\ localement
de type fini). Alors $q$ est localement de présentation finie
(resp.\ localement de type fini).
\end{lemma}

\begin{proof}
Puisque $\{X \to Y\}$ est un recouvrement fppf, il induit une surjection de
faisceaux fppf (Topologies sur les espaces, lemme
\ref{spaces-topologies-lemma-fppf-covering-surjective}); le
lemme est donc un cas particulier du lemme \ref{lemma-curiosity}.
On peut aussi le déduire plus directement de
l'analogue pour les schémas. La question est en effet locale pour la topologie \'etale
sur $B$ et sur $Y$ (Morphismes d'espaces, lemmes
\ref{spaces-morphisms-lemma-finite-type-local} et
\ref{spaces-morphisms-lemma-finite-presentation-local}).
Nous pouvons donc supposer que $B$ et $Y$ sont des
schémas affines. Comme l'application $|X| \to |Y|$ est ouverte
(Morphismes d'espaces, lemme \ref{spaces-morphisms-lemma-fppf-open}),
nous pouvons choisir un
schéma affine $U$ et un morphisme \'etale $U \to X$ tel que le
morphisme composé $U \to Y$ soit surjectif. Dans ce cas, le résultat
se déduit de Descente, lemme
\ref{descent-lemma-flat-finitely-presented-permanence}.
\end{proof}

\begin{lemma}
\label{lemma-syntomic-smooth-etale-permanence}
Soit $S$ un schéma. Considérons le diagramme suivant
$$
\xymatrix{
X \ar[rr]_f \ar[rd]_p & &
Y \ar[dl]^q \\
& B
}
$$
de morphismes d'espaces algébriques sur $S$, supposé commutatif.
Supposons que
\begin{enumerate}
\item $f$ soit surjectif et syntomique (resp.\ lisse, resp.\ \'etale),
\item $p$ soit syntomique (resp.\ lisse, resp.\ \'etale).
\end{enumerate}
Alors $q$ est syntomique (resp.\ lisse, resp.\ \'etale).
\end{lemma}

\begin{proof}
Le résultat se déduit de l'analogue pour les schémas.
La question est en effet locale pour la topologie \'etale sur $B$ et sur $Y$
(Morphismes d'espaces, lemmes
\ref{spaces-morphisms-lemma-syntomic-local},
\ref{spaces-morphisms-lemma-smooth-local}, et
\ref{spaces-morphisms-lemma-etale-local}).
Nous pouvons donc supposer que $B$ et $Y$ sont des
schémas affines. Comme l'application $|X| \to |Y|$ est ouverte
(Morphismes d'espaces, lemme \ref{spaces-morphisms-lemma-fppf-open}),
nous pouvons choisir un
schéma affine $U$ et un morphisme \'etale $U \to X$ tel que le
morphisme composé $U \to Y$ soit surjectif. Dans ce cas, le résultat
se déduit de Descente, lemme
\ref{descent-lemma-syntomic-smooth-etale-permanence}.
\end{proof}

\noindent
En fait, ce résultat peut être renforcé comme suit.

\begin{lemma}
\label{lemma-smooth-permanence}
Soit $S$ un schéma. Considérons le diagramme suivant
$$
\xymatrix{
X \ar[rr]_f \ar[rd]_p & &
Y \ar[dl]^q \\
& B
}
$$
de morphismes d'espaces algébriques sur $S$, supposé commutatif. Supposons que
\begin{enumerate}
\item $f$ soit surjectif, plat et localement de présentation finie,
\item $p$ soit lisse (resp.\ \'etale).
\end{enumerate}
Alors $q$ est lisse (resp.\ \'etale).
\end{lemma}

\begin{proof}
Le résultat se déduit de l'analogue pour les schémas.
La question est en effet locale pour la topologie \'etale sur $B$ et sur $Y$
(Morphismes d'espaces, lemmes
\ref{spaces-morphisms-lemma-smooth-local} et
\ref{spaces-morphisms-lemma-etale-local}).
Nous pouvons donc supposer que $B$ et $Y$ sont des
schémas affines. Comme l'application $|X| \to |Y|$ est ouverte
(Morphismes d'espaces, lemme \ref{spaces-morphisms-lemma-fppf-open}),
nous pouvons choisir un
schéma affine $U$ et un morphisme \'etale $U \to X$ tel que le
morphisme composé $U \to Y$ soit surjectif. Dans ce cas, le résultat
se déduit de Descente, lemme
\ref{descent-lemma-smooth-permanence}.
\end{proof}

\begin{lemma}
\label{lemma-syntomic-permanence}
Soit $S$ un schéma. Considérons le diagramme suivant
$$
\xymatrix{
X \ar[rr]_f \ar[rd]_p & &
Y \ar[dl]^q \\
& B
}
$$
de morphismes d'espaces algébriques sur $S$, supposé commutatif. Supposons que
\begin{enumerate}
\item $f$ soit surjectif, plat et localement de présentation finie,
\item $p$ soit syntomique.
\end{enumerate}
Alors $q$ et $f$ sont tous deux syntomiques.
\end{lemma}

\begin{proof}
Le résultat se déduit de l'analogue pour les schémas.
La question est en effet locale pour la topologie \'etale sur $B$ et sur $Y$
(Morphismes d'espaces, lemme
\ref{spaces-morphisms-lemma-syntomic-local}).
Nous pouvons donc supposer que $B$ et $Y$ sont des
schémas affines. Comme l'application $|X| \to |Y|$ est ouverte
(Morphismes d'espaces, lemme \ref{spaces-morphisms-lemma-fppf-open}),
nous pouvons choisir un
schéma affine $U$ et un morphisme \'etale $U \to X$ tel que le
morphisme composé $U \to Y$ soit surjectif. Dans ce cas, le résultat
se déduit de Descente, lemme
\ref{descent-lemma-syntomic-permanence}.
\end{proof}
% preserved blank authority line 843
% preserved blank authority line 844
% preserved blank authority line 845
% preserved blank authority line 846
% preserved blank authority line 847
% preserved blank authority line 848




\section{Descente de propriétés des espaces}
\label{section-descending-properties-spaces}

\noindent
Dans cette section, nous rassemblons quelques résultats du type suivant.

\begin{lemma}
\label{lemma-descend-unibranch}
Soit $S$ un schéma.
Soit $f : X \to Y$ un morphisme d'espaces algébriques sur $S$.
Soit $x \in |X|$.
Si $f$ est plat en $x$ et si $X$ est géométriquement unibranche en $x$, alors $Y$ est
géométriquement unibranche en $f(x)$.
\end{lemma}

\begin{proof}
Considérons l'homomorphisme d'anneaux strictement locaux
$\mathcal{O}_{Y, f(\overline{x})} \to \mathcal{O}_{X, \overline{x}}$.
Par
Morphismes d'espaces, lemme
\ref{spaces-morphisms-lemma-flat-at-point-etale-local-rings}
celui-ci est plat. Ainsi, si $\mathcal{O}_{X, \overline{x}}$ possède un unique idéal premier minimal,
il en va de même de $\mathcal{O}_{Y, f(\overline{x})}$ (par le théorème de descente des idéaux premiers, voir
Algèbre, lemme \ref{algebra-lemma-flat-going-down}).
\end{proof}

\begin{lemma}
\label{lemma-descend-reduced}
\begin{slogan}
Un morphisme plat et surjectif d'espaces algébriques dont la source est réduite
a un but réduit.
\end{slogan}
Soit $S$ un schéma.
Soit $f : X \to Y$ un morphisme d'espaces algébriques sur $S$.
Si $f$ est plat et surjectif et si $X$ est réduit, alors $Y$ est réduit.
\end{lemma}

\begin{proof}
Choisissons un schéma $V$ et un morphisme \'etale surjectif $V \to Y$.
Choisissons un schéma $U$ et un morphisme \'etale surjectif
$U \to X \times_Y V$. Comme $f$ est surjectif et plat, le morphisme de
schémas $U \to V$ est surjectif et plat. Nous nous ramenons ainsi
au cas des schémas (le caractère réduit de $X$ et de $Y$ étant défini
en termes du caractère réduit de $U$ et de $V$, voir
Propriétés des espaces,
section \ref{spaces-properties-section-types-properties}).
Le cas des schémas est traité dans
Descente, lemme \ref{descent-lemma-descend-reduced}.
\end{proof}

\begin{lemma}
\label{lemma-descend-locally-Noetherian}
Soit $f : X \to Y$ un morphisme d'espaces algébriques.
Si $f$ est localement de présentation finie, plat et surjectif et si
$X$ est localement noethérien, alors $Y$ est localement noethérien.
\end{lemma}

\begin{proof}
Choisissons un schéma $V$ et un morphisme \'etale surjectif $V \to Y$.
Choisissons un schéma $U$ et un morphisme \'etale surjectif
$U \to X \times_Y V$. Comme $f$ est surjectif, plat et localement de
présentation finie, le morphisme de schémas $U \to V$ est surjectif, plat et
localement de présentation finie. Nous nous ramenons ainsi
au cas des schémas (la propriété, pour $X$ et $Y$, d'être localement noethériens
étant définie par la même propriété pour $U$ et $V$, voir
Propriétés des espaces,
section \ref{spaces-properties-section-types-properties}).
Dans le cas des schémas, le résultat découle de
Descente, lemme \ref{descent-lemma-Noetherian-local-fppf}.
\end{proof}

\begin{lemma}
\label{lemma-descend-regular}
Soit $f : X \to Y$ un morphisme d'espaces algébriques.
Si $f$ est localement de présentation finie, plat et surjectif et si
$X$ est régulier, alors $Y$ est régulier.
\end{lemma}

\begin{proof}
Par
le lemme \ref{lemma-descend-locally-Noetherian}
nous savons que $Y$ est localement noethérien.
Choisissons un schéma $V$ et un morphisme \'etale surjectif $V \to Y$.
Il suffit de montrer que tous les anneaux locaux de $V$ sont réguliers ;
voir
Propriétés, lemme \ref{properties-lemma-characterize-regular}.
Choisissons un schéma $U$ et un morphisme \'etale surjectif
$U \to X \times_Y V$. Comme $f$ est surjectif et plat, le morphisme de schémas
$U \to V$ est surjectif et plat. Par hypothèse, $U$ est un schéma régulier
et, en particulier, tous ses anneaux locaux sont réguliers (d'après le lemme ci-dessus).
Le lemme découle donc de
Algèbre, lemme \ref{algebra-lemma-flat-under-regular}.
\end{proof}

\begin{lemma}
\label{lemma-reduced-local-smooth}
Soit $f : X \to Y$ un morphisme lisse d'espaces algébriques.
Si $Y$ est réduit, alors $X$ est réduit. Si $f$ est surjectif
et si $X$ est réduit, alors $Y$ est réduit.
\end{lemma}

\begin{proof}
Choisissons un diagramme commutatif
$$
\xymatrix{
U \ar[d] \ar[r] & V \ar[d] \\
X \ar[r] & Y
}
$$
où $U$ et $V$ sont des schémas, les flèches verticales sont surjectives et
\'etales, et $U \to X \times_Y V$ est surjectif et \'etale. Observons que $X$ est
un espace algébrique réduit si et seulement si $U$ est un schéma réduit
par notre définition des espaces algébriques réduits dans
Propriétés des espaces, section \ref{spaces-properties-section-types-properties}.
De même pour $Y$ et $V$. 
Le morphisme $U \to V$ est un morphisme lisse de schémas, voir
Morphismes d'espaces, lemme \ref{spaces-morphisms-lemma-smooth-local}.
Puisque le caractère réduit est local pour la topologie lisse dans la catégorie des
schémas (Descente, lemme \ref{descent-lemma-reduced-local-smooth}),
il s'ensuit que $U$ est réduit si $V$ est réduit.
D'autre part, si $X \to Y$ est surjectif, alors $U \to V$ est
surjectif et, dans ce cas, si $U$ est réduit, alors $V$ est réduit.
\end{proof}







\section{Descente de propriétés des morphismes}
\label{section-descending-properties-morphisms}

\noindent
Dans cette section, nous définissons ce que signifie, pour une propriété des morphismes
d'espaces algébriques, être locale sur le but pour une topologie. Comparer avec
Descente, section \ref{descent-section-descending-properties-morphisms}.

\begin{definition}
\label{definition-property-morphisms-local}
Soit $S$ un schéma.
Soit $\mathcal{P}$ une propriété des morphismes d'espaces algébriques sur $S$.
Soit $\tau \in \{fpqc, fppf, syntomic, smooth, \etale\}$.
On dit que $\mathcal{P}$ est {\it $\tau$-locale sur la base}, ou
{\it $\tau$-locale sur le but}, ou
{\it locale sur la base pour la topologie $\tau$} si, pour tout
$\tau$-recouvrement $\{Y_i \to Y\}_{i \in I}$ d'espaces algébriques
et tout morphisme d'espaces algébriques $f : X \to Y$, on
a
$$
f \text{ vérifie }\mathcal{P}
\Leftrightarrow
\text{chaque }Y_i \times_Y X \to Y_i\text{ vérifie }\mathcal{P}.
$$
\end{definition}

\noindent
Bien entendu, puisque les isomorphismes sont toujours des recouvrements,
on voit (ou l'on exige) que la propriété $\mathcal{P}$ est vérifiée par $X \to Y$
si et seulement si toute flèche $X' \to Y'$ isomorphe à $X \to Y$ la vérifie.
Si une propriété est $\tau$-locale sur le but, elle est préservée
par changement de base le long de morphismes figurant dans des $\tau$-recouvrements. Voici
un énoncé formel.

\begin{lemma}
\label{lemma-pullback-property-local-target}
Soit $S$ un schéma.
Soit $\tau \in \{fpqc, fppf, syntomic, smooth, \etale\}$.
Soit $\mathcal{P}$ une propriété des morphismes d'espaces algébriques sur $S$
qui est $\tau$-locale sur le but. Soit $f : X \to Y$ un morphisme vérifiant la propriété
$\mathcal{P}$. Pour tout morphisme $Y' \to Y$ qui est
plat, resp.\ plat et localement de présentation finie, resp.\ syntomique,
resp.\ \'etale, le changement de base
$f' : Y' \times_Y X \to Y'$ de $f$ vérifie la propriété $\mathcal{P}$.
\end{lemma}

\begin{proof}
Cela résulte de ce que l'on peut compléter $Y' \to Y$ en une famille de
morphismes qui constitue un $\tau$-recouvrement.
\end{proof}

\noindent
Une conséquence simple et souvent utilisée de ce qui précède est la suivante. Si
$f : X \to Y$ vérifie la propriété $\mathcal{P}$, laquelle est $\tau$-locale
sur le but, et si $f(X) \subset V$
pour un certain sous-espace ouvert $V \subset Y$, alors le morphisme induit
$X \to V$ vérifie encore $\mathcal{P}$. Démonstration : le changement de base de
$f$ par $V \to Y$ donne $X \to V$.

\begin{lemma}
\label{lemma-largest-open-of-the-base}
Soit $S$ un schéma.
Soit $\tau \in \{fppf, syntomic, smooth, \etale\}$.
Soit $\mathcal{P}$ une propriété des morphismes d'espaces algébriques sur $S$
qui est $\tau$-locale sur le but. Pour tout morphisme d'espaces algébriques
$f : X \to Y$ sur $S$, il existe un plus grand sous-espace ouvert
$W(f) \subset Y$ tel que la restriction $X_{W(f)} \to W(f)$ vérifie
$\mathcal{P}$. En outre,
\begin{enumerate}
\item si $g : Y' \to Y$ est un morphisme d'espaces algébriques qui est
plat et localement de présentation finie, syntomique, lisse ou \'etale,
et si le changement de base $f' : X_{Y'} \to Y'$ vérifie $\mathcal{P}$, alors
$g$ se factorise par $W(f)$,
\item si $g : Y' \to Y$ est plat et localement de présentation finie,
syntomique, lisse ou \'etale, alors $W(f') = g^{-1}(W(f))$, et
\item si $\{g_i : Y_i \to Y\}$ est un $\tau$-recouvrement, alors
$g_i^{-1}(W(f)) = W(f_i)$, où $f_i$ est le changement de base de $f$
par $Y_i \to Y$.
\end{enumerate}
\end{lemma}

\begin{proof}
Considérons la réunion $W_{set} \subset |Y|$ des images
$g(|Y'|) \subset |Y|$ des morphismes $g : Y' \to Y$ satisfaisant aux conditions suivantes :
\begin{enumerate}
\item $g$ est plat et localement de présentation finie, syntomique,
lisse ou \'etale, et
\item le changement de base $Y' \times_{g, Y} X \to Y'$ vérifie la propriété
$\mathcal{P}$.
\end{enumerate}
Comme un tel morphisme $g$ est ouvert (voir
Morphismes d'espaces, lemme \ref{spaces-morphisms-lemma-fppf-open}),
on voit que $W_{set}$ est une partie ouverte de $|Y|$. Notons $W \subset Y$
le sous-espace ouvert dont l'ensemble de points est $W_{set}$, voir
Propriétés des espaces, lemme \ref{spaces-properties-lemma-open-subspaces}.
Comme $\mathcal{P}$ est locale pour la topologie $\tau$, la restriction
$X_W \to W$ vérifie la propriété $\mathcal{P}$, car on dispose d'un recouvrement
$\{Y' \to W\}$ de $W$ tel que les changements de base vérifient $\mathcal{P}$.
Cela établit l'existence et montre que $W(f)$ satisfait à (1).
Pour établir (2), remarquons que $W(f') \supset g^{-1}(W(f))$, car
$\mathcal{P}$ est stable par changement de base le long de morphismes plats et localement de
présentation finie, syntomiques, lisses ou \'etales ; voir
le lemme \ref{lemma-pullback-property-local-target}.
Réciproquement, si $Y'' \subset Y'$ est un ouvert tel que
$X_{Y''} \to Y''$ vérifie la propriété $\mathcal{P}$, alors $Y'' \to Y$ se factorise
par $W$ en vertu de la construction, c'est-à-dire $Y'' \subset g^{-1}(W(f))$. Cela
prouve (2). L'assertion (3) découle de (2), car chaque morphisme
$Y_i \to Y$ est plat et localement de présentation finie, syntomique,
lisse ou \'etale d'après notre définition d'un $\tau$-recouvrement.
\end{proof}

\begin{lemma}
\label{lemma-descending-properties-morphisms}
Soit $S$ un schéma. Soit $\mathcal{P}$ une propriété des morphismes
d'espaces algébriques sur $S$. Supposons que
\begin{enumerate}
\item si les morphismes $X_i \to Y_i$, $i = 1, 2$, vérifient la propriété $\mathcal{P}$, alors
$X_1 \amalg X_2 \to Y_1 \amalg Y_2$ la vérifie aussi,
\item un morphisme d'espaces algébriques $f : X \to Y$ vérifie la propriété
$\mathcal{P}$ si et seulement si, pour tout schéma affine $Z$ et tout
morphisme $Z \to Y$, le changement de base $Z \times_Y X \to Z$ de $f$
vérifie la propriété $\mathcal{P}$, et
\item pour tout morphisme plat et surjectif de schémas affines
$Z' \to Z$ sur $S$ et tout morphisme $f : X \to Z$ d'un espace algébrique
vers $Z$, on a
$$
f' : Z' \times_Z X \to Z'\text{ vérifie }\mathcal{P}
\Rightarrow
f\text{ vérifie }\mathcal{P}.
$$
\end{enumerate}
Alors $\mathcal{P}$ est locale sur la base pour la topologie fpqc.
\end{lemma}

\begin{proof}
Si $\mathcal{P}$ satisfait à la condition (2), elle est automatiquement
stable par tout changement de base. L'implication directe de la
définition \ref{definition-property-morphisms-local} est donc acquise.

\medskip\noindent
Soit $\{Y_i \to Y\}_{i \in I}$ un recouvrement fpqc d'espaces algébriques sur $S$.
Soit $f : X \to Y$ un morphisme d'espaces algébriques sur $S$.
Supposons que chaque changement de base $f_i : Y_i \times_Y X \to Y_i$ vérifie la propriété
$\mathcal{P}$. Nous voulons montrer que $f$ vérifie $\mathcal{P}$.
Soit $Z$ un schéma affine et soit $Z \to Y$ un morphisme.
D'après (2), il suffit de montrer que le morphisme d'espaces algébriques
$Z \times_Y X \to Z$ vérifie $\mathcal{P}$.
Puisque $\{Y_i \to Y\}_{i \in I}$ est un recouvrement fpqc, il existe
un recouvrement fpqc standard $\{Z_j \to Z\}_{j = 1, \ldots , n}$
et des morphismes $Z_j \to Y_{i_j}$ au-dessus de $Y$, pour des indices convenables $i_j \in I$.
Puisque $f_{i_j}$ vérifie $\mathcal{P}$, on voit que
$$
Z_j \times_Y X
=
Z_j \times_{Y_{i_j}} (Y_{i_j} \times_Y X)
\longrightarrow
Z_j
$$
vérifie $\mathcal{P}$, car c'est un changement de base de $f_{i_j}$ (voir la première remarque de la
démonstration). Posons $Z' = \coprod_{j = 1, \ldots, n} Z_j$ ; alors $Z' \to Z$ est
un morphisme plat et surjectif de schémas affines sur $S$. D'après (1),
on conclut que $Z' \times_Y X \to Z'$ vérifie la propriété $\mathcal{P}$.
Comme il s'agit du changement de base du morphisme $Z \times_Y X \to Z$
par le morphisme $Z' \to Z$, on conclut que
$Z \times_Y X \to Z$ vérifie la propriété $\mathcal{P}$, comme souhaité.
\end{proof}
	
	
\section{Descente de propriétés des morphismes pour la topologie fpqc}
\sectionmark{Descente de propriétés pour la topologie fpqc}
\label{section-descending-properties-morphisms-fpqc}


\noindent
Dans cette section, nous montrons qu'un grand nombre de propriétés
de morphismes d'espaces algébriques sont locales sur la base
pour la topologie fpqc. Comparer avec
Descente, section \ref{descent-section-descending-properties-morphisms-fpqc}
pour le cas des morphismes de schémas.

\begin{lemma}
\label{lemma-descending-property-quasi-compact}
Soit $S$ un schéma.
La propriété $\mathcal{P}(f) =$``$f$ est quasi-compact''
est locale sur la base pour la topologie fpqc dans la catégorie des espaces algébriques sur $S$.
\end{lemma}

\begin{proof}
Nous appliquerons le
lemme \ref{lemma-descending-properties-morphisms}.
Ses hypothèses (1) et (2) résultent de
Morphismes des espaces,
Lemme \ref{spaces-morphisms-lemma-quasi-compact-local}.
Soit $Z' \to Z$ un morphisme plat et surjectif de schémas affines sur $S$.
Soit $f : X \to Z$ un morphisme d'espaces algébriques, et supposons
que le changement de base $f' : Z' \times_Z X \to Z'$ soit quasi-compact. Il faut
montrer que $f$ est quasi-compact. Il résulte de nouveau de
Morphismes des espaces,
Lemme \ref{spaces-morphisms-lemma-quasi-compact-local},
qu'il suffit de montrer que, pour tout schéma affine $Y$ et tout
morphisme $Y \to Z$, le produit fibré $Y \times_Z X$ est quasi-compact.
Voici le diagramme correspondant :
\begin{equation}
\label{equation-cube}
\vcenter{
\xymatrix{
Y \times_Z Z' \times_Z X \ar[dd] \ar[rr] \ar[rd] & &
Z' \times_Z X \ar'[d][dd]^{f'} \ar[rd] \\
& Y \times_Z X \ar[dd] \ar[rr] & & X \ar[dd]^f \\
Y \times_Z Z' \ar'[r][rr] \ar[rd] & & Z' \ar[rd] \\
& Y \ar[rr] & & Z
}
}
\end{equation}
Tous les carrés sont cartésiens, et le carré inférieur est formé
de schémas affines. L'hypothèse selon laquelle $f'$ est quasi-compact, jointe au
fait que $Y \times_Z Z'$ est affine, implique que
$Y \times_Z Z' \times_Z X$ est quasi-compact. Comme
$$
Y \times_Z Z' \times_Z X \longrightarrow Y \times_Z X
$$
est surjectif, puisqu'il s'obtient par changement de base à partir de $Z' \to Z$,
on en conclut que $Y \times_Z X$ est quasi-compact ; voir
Morphismes des espaces,
Lemme \ref{spaces-morphisms-lemma-surjection-from-quasi-compact}.
Cela achève la démonstration.
\end{proof}

\begin{lemma}
\label{lemma-descending-property-quasi-separated}
Soit $S$ un schéma.
La propriété $\mathcal{P}(f) =$``$f$ est quasi-séparé''
est locale sur la base pour la topologie fpqc dans la catégorie des espaces algébriques sur $S$.
\end{lemma}

\begin{proof}
Tout changement de base d'un morphisme quasi-séparé est quasi-séparé ; voir
Morphismes des espaces,
Lemme \ref{spaces-morphisms-lemma-base-change-separated}.
On obtient donc le sens direct de la
définition \ref{definition-property-morphisms-local}.

\medskip\noindent
Soit $\{Y_i \to Y\}_{i \in I}$ un recouvrement fpqc d'espaces algébriques sur $S$.
Soit $f : X \to Y$ un morphisme d'espaces algébriques sur $S$.
Supposons que chaque changement de base $X_i := Y_i \times_Y X \to Y_i$ soit quasi-séparé.
Cela signifie que chacun des morphismes
$$
\Delta_i :
X_i
\longrightarrow
X_i \times_{Y_i} X_i = Y_i \times_Y (X \times_Y X)
$$
est quasi-compact. Le changement de base d'un recouvrement fpqc est un recouvrement fpqc ; voir
Topologies sur les espaces, Lemme \ref{spaces-topologies-lemma-fpqc} ;
ainsi, $\{Y_i \times_Y (X \times_Y X) \to X \times_Y X\}$
est un recouvrement fpqc d'espaces algébriques. De plus, chaque
$\Delta_i$ est le changement de base du morphisme
$\Delta : X \to X \times_Y X$. Il résulte donc du
lemme \ref{lemma-descending-property-quasi-compact}
que $\Delta$ est quasi-compact, c'est-à-dire que $f$ est quasi-séparé.
\end{proof}

\begin{lemma}
\label{lemma-descending-property-universally-closed}
Soit $S$ un schéma.
La propriété $\mathcal{P}(f) =$``$f$ est universellement fermé''
est locale sur la base pour la topologie fpqc dans la catégorie des espaces algébriques sur $S$.
\end{lemma}

\begin{proof}
Nous appliquerons le
lemme \ref{lemma-descending-properties-morphisms}.
Ses hypothèses (1) et (2) résultent de
Morphismes des espaces,
Lemme \ref{spaces-morphisms-lemma-universally-closed-local}.
Soit $Z' \to Z$ un morphisme plat et surjectif de schémas affines sur $S$.
Soit $f : X \to Z$ un morphisme d'espaces algébriques, et supposons
que le changement de base $f' : Z' \times_Z X \to Z'$ soit universellement fermé.
Il faut montrer que $f$ est universellement fermé. Il résulte de nouveau de
Morphismes des espaces,
Lemme \ref{spaces-morphisms-lemma-universally-closed-local},
qu'il suffit de montrer que, pour tout schéma affine $Y$ et tout
morphisme $Y \to Z$, l'application $|Y \times_Z X| \to |Y|$ est fermée.
Considérons le cube (\ref{equation-cube}).
L'hypothèse selon laquelle $f'$ est universellement fermé implique que
$|Y \times_Z Z' \times_Z X| \to |Y \times_Z Z'|$ est fermée.
Comme $Y \times_Z Z' \to Y$ est quasi-compact, surjectif et plat,
puisqu'il s'obtient par changement de base à partir de $Z' \to Z$,
l'application $|Y \times_Z Z'| \to |Y|$ est submersive ; voir
Morphismes, Lemme \ref{morphisms-lemma-fpqc-quotient-topology}.
De plus, l'application
$$
|Y \times_Z Z' \times_Z X|
\longrightarrow
|Y \times_Z Z'| \times_{|Y|} |Y \times_Z X|
$$
est surjective ; voir
Propriétés des espaces, Lemme \ref{spaces-properties-lemma-points-cartesian}.
Il résulte de la topologie élémentaire que $|Y \times_Z X| \to |Y|$ est fermée.
\end{proof}

\begin{lemma}
\label{lemma-descending-property-universally-open}
Soit $S$ un schéma.
La propriété $\mathcal{P}(f) =$``$f$ est universellement ouvert''
est locale sur la base pour la topologie fpqc dans la catégorie des espaces algébriques sur $S$.
\end{lemma}

\begin{proof}
La démonstration est identique à celle du
lemme \ref{lemma-descending-property-universally-closed}.
\end{proof}

\begin{lemma}
\label{lemma-descending-property-universally-submersive}
La propriété $\mathcal{P}(f) =$``$f$ est universellement submersif''
est locale sur la base pour la topologie fpqc.
\end{lemma}

\begin{proof}
La démonstration est identique à celle du
lemme \ref{lemma-descending-property-universally-closed}.
\end{proof}

\begin{lemma}
\label{lemma-descending-property-surjective}
La propriété $\mathcal{P}(f) =$``$f$ est surjectif''
est locale sur la base pour la topologie fpqc.
\end{lemma}

\begin{proof}
Démonstration omise. (Indication : utiliser
Propriétés des espaces, Lemme \ref{spaces-properties-lemma-points-cartesian}.)
\end{proof}

\begin{lemma}
\label{lemma-descending-property-universally-injective}
La propriété $\mathcal{P}(f) =$``$f$ est universellement injectif''
est locale sur la base pour la topologie fpqc.
\end{lemma}

\begin{proof}
Nous appliquerons le
lemme \ref{lemma-descending-properties-morphisms}.
Ses hypothèses (1) et (2) résultent de
Morphismes des espaces,
Lemme \ref{spaces-morphisms-lemma-universally-closed-local}.
Soit $Z' \to Z$ un morphisme plat et surjectif de schémas affines
sur $S$, et soit $f : X \to Z$ un morphisme d'un espace algébrique vers $Z$.
Supposons que le changement de base $f' : X' \to Z'$ soit universellement injectif.
Soit $K$ un corps, et soient $a, b : \Spec(K) \to X$
deux morphismes tels que $f \circ a = f \circ b$.
Comme $Z' \to Z$ est surjectif, il existe une extension de corps
$K'/K$ et un morphisme
$\Spec(K') \to Z'$
de sorte que le diagramme suivant, où les flèches pleines sont données, soit commutatif :
$$
\xymatrix{
\Spec(K') \ar[rrd] \ar@{-->}[rd]_{a', b'} \ar[dd] \\
 &
X' \ar[r] \ar[d] &
Z' \ar[d] \\
\Spec(K) \ar[r]^{a, b} &
X \ar[r] &
Z
}
$$
Comme le carré est cartésien, on obtient les deux flèches pointillées $a'$, $b'$ qui rendent le
diagramme commutatif. Puisque $X' \to Z'$ est universellement injectif, on a $a' = b'$.
Cela entraîne $a = b$, car $\{\Spec(K') \to \Spec(K)\}$
est un recouvrement fpqc ; voir
Propriétés des espaces, Proposition
\ref{spaces-properties-proposition-sheaf-fpqc}.
Ainsi, $f$ est universellement injectif, ce qu'il fallait démontrer.
\end{proof}

\begin{lemma}
\label{lemma-descending-property-universal-homeomorphism}
La propriété $\mathcal{P}(f) =$``$f$ est un homéomorphisme universel''
est locale sur la base pour la topologie fpqc.
\end{lemma}

\begin{proof}
On peut le démontrer exactement de la même manière que dans le
lemme \ref{lemma-descending-property-universally-closed}.
On peut aussi utiliser le fait qu'une application d'espaces topologiques est un
homéomorphisme si et seulement si elle est injective, surjective et ouverte.
Ainsi, un homéomorphisme universel équivaut à un morphisme
surjectif, universellement injectif et universellement ouvert.
Voir Morphismes des espaces, Lemme
\ref{spaces-morphisms-lemma-base-change-surjective} et
Morphismes des espaces, Définitions
\ref{spaces-morphisms-definition-universally-injective},
\ref{spaces-morphisms-definition-surjective},
\ref{spaces-morphisms-definition-open},
\ref{spaces-morphisms-definition-universal-homeomorphism}.
Le lemme résulte donc des
lemmes \ref{lemma-descending-property-surjective},
\ref{lemma-descending-property-universally-injective} et
\ref{lemma-descending-property-universally-open}.
\end{proof}

\begin{lemma}
\label{lemma-descending-property-locally-finite-type}
La propriété $\mathcal{P}(f) =$``$f$ est localement de type fini''
est locale sur la base pour la topologie fpqc.
\end{lemma}

\begin{proof}
Nous appliquerons le
lemme \ref{lemma-descending-properties-morphisms}.
Ses hypothèses (1) et (2) résultent de
Morphismes des espaces,
Lemme \ref{spaces-morphisms-lemma-finite-type-local}.
Soit $Z' \to Z$ un morphisme plat et surjectif de schémas affines sur $S$.
Soit $f : X \to Z$ un morphisme d'espaces algébriques, et supposons
que le changement de base $f' : Z' \times_Z X \to Z'$ soit localement de type fini.
Il faut montrer que $f$ est localement de type fini. Soit $U$ un schéma
et soit $U \to X$ un morphisme étale surjectif. D'après
Morphismes des espaces,
Lemme \ref{spaces-morphisms-lemma-finite-type-local},
il suffit encore de montrer que $U \to Z$ est localement de type fini.
Puisque $f'$ est localement de type fini et que $Z' \times_Z U$ est un
schéma étale sur $Z' \times_Z X$, on déduit encore du même lemme que
$Z' \times_Z U \to Z'$ est localement de type fini.
Comme $\{Z' \to Z\}$ est un recouvrement fpqc, on en conclut que
$U \to Z$ est localement de type fini d'après
Descente, Lemme \ref{descent-lemma-descending-property-locally-finite-type},
ce qu'il fallait démontrer.
\end{proof}

\begin{lemma}
\label{lemma-descending-property-locally-finite-presentation}
La propriété $\mathcal{P}(f) =$``$f$ est localement de présentation finie''
est locale sur la base pour la topologie fpqc.
\end{lemma}

\begin{proof}
Nous appliquerons le
lemme \ref{lemma-descending-properties-morphisms}.
Ses hypothèses (1) et (2) résultent de
Morphismes des espaces,
Lemme \ref{spaces-morphisms-lemma-finite-presentation-local}.
Soit $Z' \to Z$ un morphisme plat et surjectif de schémas affines sur $S$.
Soit $f : X \to Z$ un morphisme d'espaces algébriques, et supposons
que le changement de base $f' : Z' \times_Z X \to Z'$ soit localement de
présentation finie.
Il faut montrer que $f$ est localement de présentation finie. Soit $U$ un schéma
et soit $U \to X$ un morphisme étale surjectif. D'après
Morphismes des espaces,
Lemme \ref{spaces-morphisms-lemma-finite-presentation-local},
il suffit encore de montrer que $U \to Z$ est localement de présentation finie.
Puisque $f'$ est localement de présentation finie et que $Z' \times_Z U$ est un
schéma étale sur $Z' \times_Z X$, on déduit encore du même lemme que
$Z' \times_Z U \to Z'$ est localement de présentation finie.
Comme $\{Z' \to Z\}$ est un recouvrement fpqc, on en conclut que
$U \to Z$ est localement de présentation finie d'après
Descente,
Lemme \ref{descent-lemma-descending-property-locally-finite-presentation},
ce qu'il fallait démontrer.
\end{proof}

\begin{lemma}
\label{lemma-descending-property-finite-type}
La propriété $\mathcal{P}(f) =$``$f$ est de type fini''
est locale sur la base pour la topologie fpqc.
\end{lemma}

\begin{proof}
Il suffit de combiner les lemmes \ref{lemma-descending-property-quasi-compact}
et \ref{lemma-descending-property-locally-finite-type}.
\end{proof}

\begin{lemma}
\label{lemma-descending-property-finite-presentation}
La propriété $\mathcal{P}(f) =$``$f$ est de présentation finie''
est locale sur la base pour la topologie fpqc.
\end{lemma}

\begin{proof}
Il suffit de combiner les lemmes \ref{lemma-descending-property-quasi-compact},
\ref{lemma-descending-property-quasi-separated} et
\ref{lemma-descending-property-locally-finite-presentation}.
\end{proof}

\begin{lemma}
\label{lemma-descending-property-flat}
La propriété $\mathcal{P}(f) =$``$f$ est plat''
est locale sur la base pour la topologie fpqc.
\end{lemma}

\begin{proof}
Nous appliquerons le
lemme \ref{lemma-descending-properties-morphisms}.
Ses hypothèses (1) et (2) résultent de
Morphismes des espaces,
Lemme \ref{spaces-morphisms-lemma-flat-local}.
Soit $Z' \to Z$ un morphisme plat et surjectif de schémas affines sur $S$.
Soit $f : X \to Z$ un morphisme d'espaces algébriques, et supposons
que le changement de base $f' : Z' \times_Z X \to Z'$ soit plat.
Il faut montrer que $f$ est plat. Soit $U$ un schéma
et soit $U \to X$ un morphisme étale surjectif. D'après
Morphismes des espaces,
Lemme \ref{spaces-morphisms-lemma-flat-local},
il suffit encore de montrer que $U \to Z$ est plat.
Puisque $f'$ est plat et que $Z' \times_Z U$ est un
schéma étale sur $Z' \times_Z X$, on déduit encore du même lemme que
$Z' \times_Z U \to Z'$ est plat.
Comme $\{Z' \to Z\}$ est un recouvrement fpqc, on en conclut que
$U \to Z$ est plat d'après
Descente, Lemme \ref{descent-lemma-descending-property-flat},
ce qu'il fallait démontrer.
\end{proof}

\begin{lemma}
\label{lemma-descending-property-open-immersion}
La propriété $\mathcal{P}(f) =$``$f$ est une immersion ouverte''
est locale sur la base pour la topologie fpqc.
\end{lemma}

\begin{proof}
Nous appliquerons le
lemme \ref{lemma-descending-properties-morphisms}.
Ses hypothèses (1) et (2) résultent de
Morphismes des espaces,
Lemme \ref{spaces-morphisms-lemma-closed-immersion-local}.
Considérons un diagramme cartésien
$$
\xymatrix{
X' \ar[r] \ar[d] & X \ar[d] \\
Z' \ar[r] & Z
}
$$
d'espaces algébriques sur $S$
où $Z' \to Z$ est un morphisme plat et surjectif de schémas affines
et $X' \to Z'$ une immersion ouverte. Il faut montrer que $X \to Z$
est une immersion ouverte. L'ouvert $|X'| \subset |Z'|$ correspond à un
sous-schéma ouvert $U' \subset Z'$ (isomorphe à $X'$)
tel que $\text{pr}_0^{-1}(U') = \text{pr}_1^{-1}(U')$
comme sous-schémas ouverts de $Z' \times_Z Z'$. Il existe donc un
sous-schéma ouvert $U \subset Z$ tel que $X' = (Z' \to Z)^{-1}(U)$ ; voir
Descente, Lemme \ref{descent-lemma-open-fpqc-covering}.
D'après Propriétés des espaces,
Proposition \ref{spaces-properties-proposition-sheaf-fpqc},
on voit que $X$ vérifie la condition de faisceau pour la topologie fpqc.
Nous avons maintenant le recouvrement fpqc $\mathcal{U} = \{U' \to U\}$
et l'élément $U' \to X' \to X \in \check{H}^0(\mathcal{U}, X)$.
La condition de faisceau fournit un morphisme $U \to X$ tel que
$$
\xymatrix{
U' \ar[r] \ar[d]^{\cong} \ar@/_3ex/[dd] & U \ar[d] \ar@/^3ex/[dd] \\
X' \ar[r] \ar[d] & X \ar[d] \\
Z' \ar[r] & Z
}
$$
soit commutatif. D'autre part, pour tout schéma $T$ sur $S$
et tout point à valeurs dans $T$, c'est-à-dire tout morphisme $T \to X$, le composé $T \to X \to Z$ est un
morphisme tel que $Z' \times_Z T \to Z'$ se factorise par $U'$. Cela
signifie que $T \to Z$ se factorise par $U$. Autrement dit, le morphisme
de faisceaux $U \to X$ est bijectif, ce qui achève la démonstration.
\end{proof}

\begin{lemma}
\label{lemma-descending-property-isomorphism}
La propriété $\mathcal{P}(f) =$``$f$ est un isomorphisme''
est locale sur la base pour la topologie fpqc.
\end{lemma}

\begin{proof}
Il suffit de combiner les lemmes \ref{lemma-descending-property-surjective}
et \ref{lemma-descending-property-open-immersion}.
\end{proof}

\begin{lemma}
\label{lemma-descending-property-affine}
La propriété $\mathcal{P}(f) =$``$f$ est affine''
est locale sur la base pour la topologie fpqc.
\end{lemma}

\begin{proof}
Nous appliquerons le
lemme \ref{lemma-descending-properties-morphisms}.
Ses hypothèses (1) et (2) résultent de
Morphismes des espaces,
Lemme \ref{spaces-morphisms-lemma-affine-local}.
Soit $Z' \to Z$ un morphisme plat et surjectif de schémas affines sur $S$.
Soit $f : X \to Z$ un morphisme d'espaces algébriques, et supposons
que le changement de base $f' : Z' \times_Z X \to Z'$ soit affine.
Soit $X'$ un schéma représentant $Z' \times_Z X$.
On obtient un isomorphisme canonique
$$
\varphi : X' \times_Z Z' \longrightarrow Z' \times_Z X'
$$
car les deux schémas représentent l'espace algébrique $Z' \times_Z Z' \times_Z X$.
C'est une donnée de descente pour $X'/Z'/Z$ ; voir
Descente, Définition \ref{descent-definition-descent-datum}
(vérification omise ; comparer avec
Descente, Lemme \ref{descent-lemma-descent-data-sheaves}).
Puisque $X' \to Z'$ est affine, cette donnée de descente est effective ; voir
Descente, Lemme \ref{descent-lemma-affine}.
Il existe donc un schéma $Y \to Z$ sur $Z$, ainsi qu'un
isomorphisme $\psi : Z' \times_Z Y \to X'$ compatible avec les données de descente.
Bien entendu, $Y \to Z$ est affine (par construction ou d'après
Descente, Lemme \ref{descent-lemma-descending-property-affine}).
Remarquons que $\mathcal{Y} = \{Z' \times_Z Y \to Y\}$ est un
recouvrement fpqc. En interprétant $\psi$ comme un élément de
$X(Z' \times_Z Y)$, on voit que $\psi \in \check{H}^0(\mathcal{Y}, X)$.
La condition de faisceau pour $X$ relativement à ce recouvrement (voir
Propriétés des espaces, Proposition
\ref{spaces-properties-proposition-sheaf-fpqc})
fournit un morphisme $Y \to X$.
Par construction, le changement de base de ce morphisme par $Z'$ est un isomorphisme ; le morphisme lui-même est donc
un isomorphisme d'après le
lemme \ref{lemma-descending-property-isomorphism}.
Cela prouve que $X$ est représentable par un schéma affine, ce qui achève la démonstration.
\end{proof}

\begin{lemma}
\label{lemma-descending-property-closed-immersion}
La propriété $\mathcal{P}(f) =$``$f$ est une immersion fermée''
est locale sur la base pour la topologie fpqc.
\end{lemma}

\begin{proof}
Nous appliquerons le
lemme \ref{lemma-descending-properties-morphisms}.
Ses hypothèses (1) et (2) résultent de
Morphismes des espaces,
Lemme \ref{spaces-morphisms-lemma-closed-immersion-local}.
Considérons un diagramme cartésien
$$
\xymatrix{
X' \ar[r] \ar[d] & X \ar[d] \\
Z' \ar[r] & Z
}
$$
d'espaces algébriques sur $S$
où $Z' \to Z$ est un morphisme plat et surjectif de schémas affines
et $X' \to Z'$ une immersion fermée. Il faut montrer que $X \to Z$
est une immersion fermée. Le morphisme $X' \to Z'$ est affine. D'après le
lemme \ref{lemma-descending-property-affine},
on voit que $X$ est un schéma et que $X \to Z$ est affine.
Il résulte de
Descente, Lemme \ref{descent-lemma-descending-property-closed-immersion},
que $X \to Z$ est une immersion fermée, ce qu'il fallait démontrer.
\end{proof}

\begin{lemma}
\label{lemma-descending-property-separated}
La propriété $\mathcal{P}(f) =$``$f$ est séparé''
est locale sur la base pour la topologie fpqc.
\end{lemma}

\begin{proof}
Tout changement de base d'un morphisme séparé est séparé ; voir
Morphismes des espaces,
Lemme \ref{spaces-morphisms-lemma-base-change-separated}.
Cela donne donc l'implication directe de la
définition \ref{definition-property-morphisms-local}.

\medskip\noindent
Soit $\{Y_i \to Y\}_{i \in I}$ un recouvrement fpqc d'espaces algébriques sur $S$.
Soit $f : X \to Y$ un morphisme d'espaces algébriques sur $S$.
Supposons que chaque changement de base $X_i := Y_i \times_Y X \to Y_i$ soit séparé.
Cela signifie que chacun des morphismes
$$
\Delta_i :
X_i
\longrightarrow
X_i \times_{Y_i} X_i = Y_i \times_Y (X \times_Y X)
$$
est une immersion fermée. Le changement de base d'un recouvrement fpqc est un
recouvrement fpqc ; voir
Topologies sur les espaces, Lemme \ref{spaces-topologies-lemma-fpqc} ;
ainsi, $\{Y_i \times_Y (X \times_Y X) \to X \times_Y X\}$
est un recouvrement fpqc d'espaces algébriques. De plus, chaque
$\Delta_i$ est le changement de base du morphisme
$\Delta : X \to X \times_Y X$. Il résulte donc du
lemme \ref{lemma-descending-property-closed-immersion}
que $\Delta$ est une immersion fermée, c'est-à-dire que $f$ est séparé.
\end{proof}

\begin{lemma}
\label{lemma-descending-property-proper}
La propriété $\mathcal{P}(f) =$``$f$ est propre''
est locale sur la base pour la topologie fpqc.
\end{lemma}

\begin{proof}
Le lemme résulte de la combinaison des
lemmes \ref{lemma-descending-property-universally-closed},
\ref{lemma-descending-property-separated}
et \ref{lemma-descending-property-finite-type}.
\end{proof}

\begin{lemma}
\label{lemma-descending-property-quasi-affine}
La propriété $\mathcal{P}(f) =$``$f$ est quasi-affine''
est locale sur la base pour la topologie fpqc.
\end{lemma}

\begin{proof}
Nous appliquerons le
lemme \ref{lemma-descending-properties-morphisms}.
Ses hypothèses (1) et (2) résultent de
Morphismes des espaces,
Lemme \ref{spaces-morphisms-lemma-quasi-affine-local}.
Soit $Z' \to Z$ un morphisme plat et surjectif de schémas affines sur $S$.
Soit $f : X \to Z$ un morphisme d'espaces algébriques, et supposons
que le changement de base $f' : Z' \times_Z X \to Z'$ soit quasi-affine.
Soit $X'$ un schéma représentant $Z' \times_Z X$.
On obtient un isomorphisme canonique
$$
\varphi : X' \times_Z Z' \longrightarrow Z' \times_Z X'
$$
car les deux schémas représentent l'espace algébrique $Z' \times_Z Z' \times_Z X$.
C'est une donnée de descente pour $X'/Z'/Z$ ; voir
Descente, Définition \ref{descent-definition-descent-datum}
(vérification omise ; comparer avec
Descente, Lemme \ref{descent-lemma-descent-data-sheaves}).
Puisque $X' \to Z'$ est quasi-affine, cette donnée de descente est effective ; voir
Descente, Lemme \ref{descent-lemma-quasi-affine}.
Il existe donc un schéma $Y \to Z$ sur $Z$, ainsi qu'un
isomorphisme $\psi : Z' \times_Z Y \to X'$ compatible avec les données de descente.
Bien entendu, $Y \to Z$ est quasi-affine (par construction ou d'après
Descente, Lemme \ref{descent-lemma-descending-property-quasi-affine}).
Remarquons que $\mathcal{Y} = \{Z' \times_Z Y \to Y\}$ est un
recouvrement fpqc. En interprétant $\psi$ comme un élément de
$X(Z' \times_Z Y)$, on voit que $\psi \in \check{H}^0(\mathcal{Y}, X)$.
La condition de faisceau pour $X$ (voir
Propriétés des espaces, Proposition
\ref{spaces-properties-proposition-sheaf-fpqc})
fournit un morphisme $Y \to X$.
Par construction, le changement de base de ce morphisme par $Z'$ est un isomorphisme ; il est donc
un isomorphisme d'après le
lemme \ref{lemma-descending-property-isomorphism}.
Cela prouve que $X$ est représentable par un schéma quasi-affine, ce qui achève la démonstration.
\end{proof}

\begin{lemma}
\label{lemma-descending-property-quasi-compact-immersion}
La propriété $\mathcal{P}(f) =$``$f$ est une immersion quasi-compacte''
est locale sur la base pour la topologie fpqc.
\end{lemma}

\begin{proof}
Nous appliquerons le
lemme \ref{lemma-descending-properties-morphisms}.
Ses hypothèses (1) et (2) résultent de
Morphismes des espaces,
Lemmes \ref{spaces-morphisms-lemma-closed-immersion-local} et
\ref{spaces-morphisms-lemma-quasi-compact-local}.
Considérons un diagramme cartésien
$$
\xymatrix{
X' \ar[r] \ar[d] & X \ar[d] \\
Z' \ar[r] & Z
}
$$
d'espaces algébriques sur $S$
où $Z' \to Z$ est un morphisme plat et surjectif de schémas affines
et $X' \to Z'$ une immersion quasi-compacte. Il faut montrer que $X \to Z$
est une immersion quasi-compacte. Le morphisme $X' \to Z'$ est quasi-affine. D'après le
lemme \ref{lemma-descending-property-quasi-affine},
on voit que $X$ est un schéma et que $X \to Z$ est quasi-affine.
Il résulte de
Descente, Lemme \ref{descent-lemma-descending-property-quasi-compact-immersion},
que $X \to Z$ est une immersion quasi-compacte, ce qu'il fallait démontrer.
\end{proof}

\begin{lemma}
\label{lemma-descending-property-integral}
La propriété $\mathcal{P}(f) =$``$f$ est entier''
est locale sur la base pour la topologie fpqc.
\end{lemma}

\begin{proof}
Un morphisme est entier si et seulement s'il est affine et
universellement fermé. Voir
Morphismes des espaces,
Lemme \ref{spaces-morphisms-lemma-integral-universally-closed}.
Le lemme résulte donc de la combinaison des
lemmes \ref{lemma-descending-property-universally-closed}
et \ref{lemma-descending-property-affine}.
\end{proof}

\begin{lemma}
\label{lemma-descending-property-finite}
La propriété $\mathcal{P}(f) =$``$f$ est fini''
est locale sur la base pour la topologie fpqc.
\end{lemma}

\begin{proof}
Un morphisme est fini si et seulement s'il est entier
et localement de type fini. Voir
Morphismes des espaces, Lemme \ref{spaces-morphisms-lemma-finite-integral}.
Le lemme résulte donc de la combinaison des
lemmes \ref{lemma-descending-property-locally-finite-type}
et \ref{lemma-descending-property-integral}.
\end{proof}

\begin{lemma}
\label{lemma-descending-property-quasi-finite}
Les propriétés
$\mathcal{P}(f) =$``$f$ est localement quasi-fini''
et
$\mathcal{P}(f) =$``$f$ est quasi-fini''
sont locales sur la base pour la topologie fpqc.
\end{lemma}

\begin{proof}
Nous avons déjà vu que la propriété ``quasi-compact'' est locale sur la base pour la topologie fpqc ; voir
Lemme \ref{lemma-descending-property-quasi-compact}. Il suffit donc
de démontrer le lemme pour ``localement quasi-fini''. Nous appliquerons le
lemme \ref{lemma-descending-properties-morphisms}.
Ses hypothèses (1) et (2) résultent de
Morphismes des espaces,
Lemme \ref{spaces-morphisms-lemma-quasi-finite-local}.
Soit $Z' \to Z$ un morphisme plat et surjectif de schémas affines sur $S$.
Soit $f : X \to Z$ un morphisme d'espaces algébriques, et supposons
que le changement de base $f' : Z' \times_Z X \to Z'$ soit localement quasi-fini.
Il faut montrer que $f$ est localement quasi-fini. Soit $U$ un schéma
et soit $U \to X$ un morphisme étale surjectif. D'après
Morphismes des espaces,
Lemme \ref{spaces-morphisms-lemma-quasi-finite-local},
il suffit encore de montrer que $U \to Z$ est localement quasi-fini.
Puisque $f'$ est localement quasi-fini et que $Z' \times_Z U$ est un
schéma étale sur $Z' \times_Z X$, on déduit encore du même lemme que
$Z' \times_Z U \to Z'$ est localement quasi-fini.
Comme $\{Z' \to Z\}$ est un recouvrement fpqc, on en conclut que
$U \to Z$ est localement quasi-fini d'après
Descente, Lemme \ref{descent-lemma-descending-property-quasi-finite},
ce qu'il fallait démontrer.
\end{proof}

\begin{lemma}
\label{lemma-descending-property-syntomic}
La propriété $\mathcal{P}(f) =$``$f$ est syntomique''
est locale sur la base pour la topologie fpqc.
\end{lemma}

\begin{proof}
Nous appliquerons le
lemme \ref{lemma-descending-properties-morphisms}.
Ses hypothèses (1) et (2) résultent de
Morphismes des espaces,
Lemme \ref{spaces-morphisms-lemma-syntomic-local}.
Soit $Z' \to Z$ un morphisme plat et surjectif de schémas affines sur $S$.
Soit $f : X \to Z$ un morphisme d'espaces algébriques, et supposons
que le changement de base $f' : Z' \times_Z X \to Z'$ soit syntomique.
Il faut montrer que $f$ est syntomique. Soit $U$ un schéma
et soit $U \to X$ un morphisme étale surjectif. D'après
Morphismes des espaces,
Lemme \ref{spaces-morphisms-lemma-syntomic-local},
il suffit encore de montrer que $U \to Z$ est syntomique.
Puisque $f'$ est syntomique et que $Z' \times_Z U$ est un
schéma étale sur $Z' \times_Z X$, on déduit encore du même lemme que
$Z' \times_Z U \to Z'$ est syntomique.
Comme $\{Z' \to Z\}$ est un recouvrement fpqc, on en conclut que
$U \to Z$ est syntomique d'après
Descente, Lemme \ref{descent-lemma-descending-property-syntomic},
ce qu'il fallait démontrer.
\end{proof}

\begin{lemma}
\label{lemma-descending-property-smooth}
La propriété $\mathcal{P}(f) =$``$f$ est lisse''
est locale sur la base pour la topologie fpqc.
\end{lemma}

\begin{proof}
Nous appliquerons le
lemme \ref{lemma-descending-properties-morphisms}.
Ses hypothèses (1) et (2) résultent de
Morphismes des espaces,
Lemme \ref{spaces-morphisms-lemma-smooth-local}.
Soit $Z' \to Z$ un morphisme plat et surjectif de schémas affines sur $S$.
Soit $f : X \to Z$ un morphisme d'espaces algébriques, et supposons
que le changement de base $f' : Z' \times_Z X \to Z'$ soit lisse.
Il faut montrer que $f$ est lisse. Soit $U$ un schéma
et soit $U \to X$ un morphisme étale surjectif. D'après
Morphismes des espaces,
Lemme \ref{spaces-morphisms-lemma-smooth-local},
il suffit encore de montrer que $U \to Z$ est lisse.
Puisque $f'$ est lisse et que $Z' \times_Z U$ est un
schéma étale sur $Z' \times_Z X$, on déduit encore du même lemme que
$Z' \times_Z U \to Z'$ est lisse.
Comme $\{Z' \to Z\}$ est un recouvrement fpqc, on en conclut que
$U \to Z$ est lisse d'après
Descente, Lemme \ref{descent-lemma-descending-property-smooth},
ce qu'il fallait démontrer.
\end{proof}

\begin{lemma}
\label{lemma-descending-property-unramified}
La propriété $\mathcal{P}(f) =$``$f$ est non ramifié''
est locale sur la base pour la topologie fpqc.
\end{lemma}

\begin{proof}
Nous appliquerons le
lemme \ref{lemma-descending-properties-morphisms}.
Ses hypothèses (1) et (2) résultent de
Morphismes des espaces,
Lemme \ref{spaces-morphisms-lemma-unramified-local}.
Soit $Z' \to Z$ un morphisme plat et surjectif de schémas affines sur $S$.
Soit $f : X \to Z$ un morphisme d'espaces algébriques, et supposons
que le changement de base $f' : Z' \times_Z X \to Z'$ soit non ramifié.
Il faut montrer que $f$ est non ramifié. Soit $U$ un schéma
et soit $U \to X$ un morphisme étale surjectif. D'après
Morphismes des espaces,
Lemme \ref{spaces-morphisms-lemma-unramified-local},
il suffit encore de montrer que $U \to Z$ est non ramifié.
Puisque $f'$ est non ramifié et que $Z' \times_Z U$ est un
schéma étale sur $Z' \times_Z X$, on déduit encore du même lemme que
$Z' \times_Z U \to Z'$ est non ramifié.
Comme $\{Z' \to Z\}$ est un recouvrement fpqc, on en conclut que
$U \to Z$ est non ramifié d'après
Descente, Lemme \ref{descent-lemma-descending-property-unramified},
ce qu'il fallait démontrer.
\end{proof}

\begin{lemma}
\label{lemma-descending-property-etale}
La propriété $\mathcal{P}(f) =$``$f$ est étale''
est locale sur la base pour la topologie fpqc.
\end{lemma}

\begin{proof}
Nous appliquerons le
lemme \ref{lemma-descending-properties-morphisms}.
Ses hypothèses (1) et (2) résultent de
Morphismes des espaces,
Lemme \ref{spaces-morphisms-lemma-etale-local}.
Soit $Z' \to Z$ un morphisme plat et surjectif de schémas affines sur $S$.
Soit $f : X \to Z$ un morphisme d'espaces algébriques, et supposons
que le changement de base $f' : Z' \times_Z X \to Z'$ soit étale.
Il faut montrer que $f$ est étale. Soit $U$ un schéma
et soit $U \to X$ un morphisme étale surjectif. D'après
Morphismes des espaces,
Lemme \ref{spaces-morphisms-lemma-etale-local},
il suffit encore de montrer que $U \to Z$ est étale.
Puisque $f'$ est étale et que $Z' \times_Z U$ est un
schéma étale sur $Z' \times_Z X$, on déduit encore du même lemme que
$Z' \times_Z U \to Z'$ est étale.
Comme $\{Z' \to Z\}$ est un recouvrement fpqc, on en conclut que
$U \to Z$ est étale d'après
Descente, Lemme \ref{descent-lemma-descending-property-etale},
ce qu'il fallait démontrer.
\end{proof}

\begin{lemma}
\label{lemma-descending-property-finite-locally-free}
La propriété $\mathcal{P}(f) =$``$f$ est fini localement libre''
est locale sur la base pour la topologie fpqc.
\end{lemma}

\begin{proof}
Un morphisme est fini localement libre si et seulement s'il est
fini, plat et localement de présentation finie
(Morphismes des espaces, Lemme \ref{spaces-morphisms-lemma-finite-flat}).
Le résultat découle donc des lemmes
\ref{lemma-descending-property-finite},
\ref{lemma-descending-property-flat} et
\ref{lemma-descending-property-locally-finite-presentation}.
\end{proof}

\begin{lemma}
\label{lemma-descending-property-monomorphism}
La propriété $\mathcal{P}(f) =$``$f$ est un monomorphisme''
est locale sur la base pour la topologie fpqc.
\end{lemma}

\begin{proof}
Soit $f : X \to Y$ un morphisme d'espaces algébriques.
Soit $\{Y_i \to Y\}$ un recouvrement fpqc, et supposons
que chacun des changements de base $f_i : X_i \to Y_i$ de $f$ soit
un monomorphisme. Il faut montrer que $f$ est un monomorphisme.

\medskip\noindent
Première démonstration. Le morphisme $f$ est un monomorphisme si et seulement si
$\Delta : X \to X \times_Y X$ est un isomorphisme. En appliquant ce critère aux
$f_i$, on voit que chacun des morphismes
$$
\Delta_i :
X_i
\longrightarrow
X_i \times_{Y_i} X_i = Y_i \times_Y (X \times_Y X)
$$
est un isomorphisme. Le changement de base d'un recouvrement fpqc est un recouvrement fpqc ; voir
Topologies sur les espaces, Lemme \ref{spaces-topologies-lemma-fpqc} ;
ainsi, $\{Y_i \times_Y (X \times_Y X) \to X \times_Y X\}$
est un recouvrement fpqc d'espaces algébriques. De plus, chaque
$\Delta_i$ est le changement de base du morphisme
$\Delta : X \to X \times_Y X$. Il résulte donc du
lemme \ref{lemma-descending-property-isomorphism}
que $\Delta$ est un isomorphisme, c'est-à-dire que $f$ est un monomorphisme.

\medskip\noindent
Seconde démonstration.
Soit $V$ un schéma et soit $V \to Y$ un morphisme étale surjectif.
Si l'on peut montrer que $V \times_Y X \to V$ est un monomorphisme, alors
$X \to Y$ est un monomorphisme. En effet, dans tout
diagramme cartésien de faisceaux
$$
\vcenter{
\xymatrix{
\mathcal{F} \ar[r]_a \ar[d]_b & \mathcal{G} \ar[d]^c \\
\mathcal{H} \ar[r]^d & \mathcal{I}
}
}
\quad
\quad
\mathcal{F} = \mathcal{H} \times_\mathcal{I} \mathcal{G}
$$
si $c$ est un morphisme surjectif de faisceaux et si $a$ est injectif, alors
$d$ est lui aussi injectif. Cela ramène le problème au cas où $Y$ est
un schéma. En outre, dans ce cas, on peut supposer que les espaces algébriques
$Y_i$ sont des schémas, puisqu'on peut toujours raffiner le recouvrement afin de se
ramener à cette situation ; voir
Topologies sur les espaces, Lemme \ref{spaces-topologies-lemma-refine-fpqc-schemes}.

\medskip\noindent
Supposons que $\{Y_i \to Y\}$ soit un recouvrement fpqc de schémas.
Soient $a, b : T \to X$ deux morphismes
tels que $f \circ a = f \circ b$. Il faut montrer que $a = b$.
Puisque $f_i$ est un monomorphisme, on a $a_i = b_i$, où
$a_i, b_i : Y_i \times_Y T \to X_i$ sont
les changements de base. En particulier, les composés
$Y_i \times_Y T \to T \to X$ sont égaux.
Comme $\{Y_i \times_Y T \to T\}$ est un recouvrement fpqc,
on déduit que $a = b$ d'après Propriétés des espaces, Proposition
\ref{spaces-properties-proposition-sheaf-fpqc}.
\end{proof}



% preserved blank authority line 2020
\section{Descente des propriétés des morphismes pour la topologie fppf}
\label{section-descending-properties-morphisms-fppf}

\noindent
Nous établissons dans cette section certaines propriétés des morphismes d'espaces
algébriques dont nous n'avons pas (encore) pu montrer qu'elles sont locales sur
la base pour la topologie fpqc, mais qui sont locales sur la base
pour la topologie fppf.

\begin{lemma}
\label{lemma-descending-fppf-property-immersion}
La propriété $\mathcal{P}(f) =$``$f$ est une immersion''
est locale sur la base pour la topologie fppf.
\end{lemma}

\begin{proof}
Soit $f : X \to Y$ un morphisme d'espaces algébriques.
Soit $\{Y_i \to Y\}_{i \in I}$ un recouvrement fppf de $Y$.
Soit $f_i : X_i \to Y_i$ le changement de base de $f$.

\medskip\noindent
Si $f$ est une immersion, alors chaque $f_i$ est une immersion d'après
Espaces, lemme \ref{spaces-lemma-base-change-immersions}.
Cela démontre l'implication directe de la
définition \ref{definition-property-morphisms-local}.

\medskip\noindent
Réciproquement, supposons que chaque $f_i$ soit une immersion. D'après
Morphismes d'espaces,
lemme \ref{spaces-morphisms-lemma-immersions-monomorphisms},
chaque $f_i$ est alors séparé. D'après
Morphismes d'espaces,
lemme \ref{spaces-morphisms-lemma-immersion-quasi-finite},
chaque $f_i$ est alors localement quasi-fini.
Les lemmes \ref{lemma-descending-property-separated}
et \ref{lemma-descending-property-quasi-finite} montrent donc que
$f$ est localement quasi-fini et séparé.
D'après
Morphismes d'espaces, lemme
\ref{spaces-morphisms-lemma-locally-quasi-finite-separated-representable},
il s'ensuit que $f$ est représentable !

\medskip\noindent
D'après
Morphismes d'espaces, lemme \ref{spaces-morphisms-lemma-closed-immersion-local},
il suffit de montrer que, pour tout schéma $Z$ et tout morphisme $Z \to Y$,
le changement de base $Z \times_Y X \to Z$ est une immersion. D'après
Topologies sur les espaces, lemme \ref{spaces-topologies-lemma-refine-fppf-schemes},
on peut trouver un recouvrement fppf $\{Z_i \to Z\}$ par des schémas qui raffine
le changement de base du recouvrement $\{Y_i \to Y\}$ à $Z$.
Par conséquent, $Z \times_Y X \to Z$ (qui est un morphisme de schémas
d'après le résultat du paragraphe précédent) devient une immersion
après changement de base par chacun des morphismes d'un recouvrement fppf de $Z$ formé de schémas.
Ainsi, $Z \times_Y X \to Z$ est une immersion d'après le résultat pour les schémas ; voir
Descente, lemme \ref{descent-lemma-descending-fppf-property-immersion}.
\end{proof}

\begin{lemma}
\label{lemma-descending-fppf-property-locally-separated}
La propriété $\mathcal{P}(f) =$``$f$ est localement séparé''
est locale sur la base pour la topologie fppf.
\end{lemma}

\begin{proof}
Tout changement de base d'un morphisme localement séparé est localement séparé ; voir
Morphismes d'espaces,
lemme \ref{spaces-morphisms-lemma-base-change-separated}.
Cela établit l'implication directe de la
définition \ref{definition-property-morphisms-local}.

\medskip\noindent
Soit $\{Y_i \to Y\}_{i \in I}$ un recouvrement fppf d'espaces algébriques sur $S$.
Soit $f : X \to Y$ un morphisme d'espaces algébriques sur $S$.
Supposons que chaque changement de base $X_i := Y_i \times_Y X \to Y_i$ soit localement séparé.
Cela signifie que chacun des morphismes
$$
\Delta_i :
X_i
\longrightarrow
X_i \times_{Y_i} X_i = Y_i \times_Y (X \times_Y X)
$$
est une immersion. Tout changement de base d'un recouvrement fppf est un
recouvrement fppf ; voir
Topologies sur les espaces, lemme \ref{spaces-topologies-lemma-fppf}.
Par conséquent, $\{Y_i \times_Y (X \times_Y X) \to X \times_Y X\}$
est un recouvrement fppf d'espaces algébriques. De plus, chaque
$\Delta_i$ est le changement de base du morphisme
$\Delta : X \to X \times_Y X$. Il résulte donc du
lemme \ref{lemma-descending-fppf-property-immersion}
que $\Delta$ est une immersion, c'est-à-dire que $f$ est localement séparé.
\end{proof}





\section{Application de la descente des propriétés des morphismes}
\label{section-application-descending-properties-morphisms}

\noindent
Cette section est l'analogue de Descente, section
\ref{descent-section-application-descending-properties-morphisms}.

\begin{lemma}
\label{lemma-descending-property-ample}
Soit $S$ un schéma.
Soit $f : X \to Y$ un morphisme d'espaces algébriques sur $S$.
Soit $\mathcal{L}$ un $\mathcal{O}_X$-module inversible.
Soit $\{g_i : Y_i \to Y\}_{i \in I}$ un recouvrement fpqc.
Soit $f_i : X_i \to Y_i$ le changement de base de $f$, et soit $\mathcal{L}_i$
l'image inverse de $\mathcal{L}$ sur $X_i$.
Les conditions suivantes sont équivalentes :
\begin{enumerate}
\item $\mathcal{L}$ est ample sur $X/Y$ ; et
\item $\mathcal{L}_i$ est ample sur $X_i/Y_i$
pour tout $i \in I$.
\end{enumerate}
\end{lemma}

\begin{proof}
L'implication (1) $\Rightarrow$ (2) résulte de
Diviseurs sur les espaces, lemme \ref{spaces-divisors-lemma-ample-base-change}.
Supposons (2). Pour vérifier que $\mathcal{L}$ est ample sur $X/Y$, on peut
travailler localement pour la topologie \'etale sur $Y$ ; voir
Diviseurs sur les espaces, lemme \ref{spaces-divisors-lemma-relatively-ample-local}.
On peut donc supposer que $Y$ est un schéma, puis que
chaque $Y_i$ est également un schéma ; voir
Topologies sur les espaces, lemme \ref{spaces-topologies-lemma-refine-fpqc-schemes}.
Autrement dit, on peut supposer que
$\{Y_i \to Y\}$ est un recouvrement fpqc de schémas.

\medskip\noindent
D'après Diviseurs sur les espaces, lemme
\ref{spaces-divisors-lemma-relatively-ample-properties},
$X_i \to Y_i$ est représentable (c'est-à-dire que $X_i$ est un schéma),
quasi-compact et séparé. Les lemmes
\ref{lemma-descending-property-quasi-compact} et
\ref{lemma-descending-property-separated} montrent donc que $f$ est quasi-compact et séparé.
Cela signifie que
$\mathcal{A} = \bigoplus_{d \geq 0} f_*\mathcal{L}^{\otimes d}$
est une algèbre graduée quasi-cohérente sur $\mathcal{O}_Y$
(Morphismes d'espaces, lemme \ref{spaces-morphisms-lemma-pushforward}).
De plus, la formation de $\mathcal{A}$ commute aux
changements de base plats d'après
Cohomologie des espaces, lemme
\ref{spaces-cohomology-lemma-flat-base-change-cohomology}.
En particulier, si l'on pose
$\mathcal{A}_i = \bigoplus_{d \geq 0} f_{i, *}\mathcal{L}_i^{\otimes d}$,
alors $\mathcal{A}_i = g_i^*\mathcal{A}$.
Il s'ensuit que les morphismes naturels
$\psi_d : f^*\mathcal{A}_d \to \mathcal{L}^{\otimes d}$
de $\mathcal{O}_X$-modules
deviennent, après image inverse, les morphismes naturels
$\psi_{i, d} : f_i^*(\mathcal{A}_i)_d \to \mathcal{L}_i^{\otimes d}$
de $\mathcal{O}_{X_i}$-modules. Puisque $\mathcal{L}_i$ est ample sur $X_i/Y_i$,
pour tout point $x_i \in X_i$, il existe un $d \geq 1$
tel que $f_i^*(\mathcal{A}_i)_d \to \mathcal{L}_i^{\otimes d}$
soit surjectif sur la fibre en $x_i$. Cela résulte soit directement
de la définition d'un module relativement ample, soit de
Morphismes, lemme \ref{morphisms-lemma-characterize-relatively-ample}.
Si $x \in |X|$, on peut choisir un $i$ et un $x_i \in X_i$
dont l'image est $x$. Comme
$\mathcal{O}_{X, \overline{x}} \to \mathcal{O}_{X_i, \overline{x}_i}$
est plat, donc fidèlement plat, on conclut que, pour tout $x \in |X|$,
il existe un $d \geq 1$ tel que
$f^*\mathcal{A}_d \to \mathcal{L}^{\otimes d}$
soit surjectif sur la fibre en $x$.
Il s'ensuit que l'ouvert $U(\psi) \subset X$ du
lemme de Diviseurs sur les espaces
\ref{spaces-divisors-lemma-invertible-map-into-relative-proj}
correspondant au morphisme
$\psi : f^*\mathcal{A} \to \bigoplus_{d \geq 0} \mathcal{L}^{\otimes d}$
d'algèbres graduées sur $\mathcal{O}_X$
est égal à $X$. Considérons le morphisme correspondant
$$
r_{\mathcal{L}, \psi} : X \longrightarrow \underline{\text{Proj}}_Y(\mathcal{A})
$$
Il ressort clairement de ce qui précède que le changement de base de
$r_{\mathcal{L}, \psi}$ à $Y_i$ est le morphisme
$r_{\mathcal{L}_i, \psi_i}$, qui est une immersion ouverte d'après
Morphismes, lemme \ref{morphisms-lemma-characterize-relatively-ample}.
Ainsi, $r_{\mathcal{L}, \psi}$ est une immersion ouverte
d'après le lemme \ref{lemma-descending-property-open-immersion}.
Par conséquent, $X$ est un schéma,
et l'on conclut que $\mathcal{L}$ est ample sur $X/Y$ d'après
Morphismes, lemme \ref{morphisms-lemma-characterize-relatively-ample}.
\end{proof}

\begin{lemma}
\label{lemma-ample-in-neighbourhood}
Soit $S$ un schéma.
Soit $f : X \to Y$ un morphisme propre d'espaces algébriques sur $S$.
Soit $\mathcal{L}$ un $\mathcal{O}_X$-module inversible.
Il existe un sous-espace ouvert $V \subset Y$ caractérisé par
la propriété suivante :
Un morphisme $Y' \to Y$ d'espaces algébriques se factorise
par $V$ si et seulement si l'image inverse $\mathcal{L}'$
de $\mathcal{L}$ sur $X' = Y' \times_Y X$ est ample sur $X'/Y'$
(au sens de Diviseurs sur les espaces, définition
\ref{spaces-divisors-definition-relatively-ample}).
\end{lemma}

\begin{proof}
Supposons le lemme établi lorsque $Y$ est un schéma.
Soit $U$ un schéma et soit $U \to Y$ un morphisme \'etale surjectif.
Soit $R = U \times_Y U$, avec les projections $t, s : R \to U$.
Notons $X_U = U \times_Y X$ et désignons par $\mathcal{L}_U$ l'image inverse correspondante.
On obtient alors un sous-schéma ouvert $V' \subset U$ comme dans le lemme
pour $(X_U \to U, \mathcal{L}_U)$. La caractérisation fonctorielle donne
$s^{-1}(V') = t^{-1}(V')$.
Il existe donc un sous-espace ouvert $V \subset Y$ tel que
$V'$ soit l'image inverse de $V$ dans $U$.
En particulier, $V' \to V$ est \'etale et surjectif, et l'on
conclut que $\mathcal{L}_V$ est ample sur $X_V/V$
(Diviseurs sur les espaces, lemme \ref{spaces-divisors-lemma-relatively-ample-local}).
Maintenant, si $Y' \to Y$ est un morphisme tel que
$\mathcal{L}'$ soit ample sur $X'/Y'$, alors
$U \times_Y Y' \to Y'$ doit se factoriser par $V'$,
et l'on conclut que $Y' \to Y$ se factorise par $V$.
Ainsi, $V \subset Y$ possède la propriété énoncée dans le lemme.
On est ainsi ramené au cas traité au
paragraphe suivant.

\medskip\noindent
Supposons que $Y$ soit un schéma. La question étant locale sur $Y$,
on peut supposer que $Y$ est un schéma affine. Montrons
l'assertion suivante :
\begin{enumerate}
\item[(A)] Si $\Spec(k) \to Y$ est un morphisme tel que
$\mathcal{L}_k$ soit ample sur $X_k/k$, alors il existe un
voisinage ouvert $V \subset Y$ de l'image de $\Spec(k) \to Y$
tel que $\mathcal{L}_V$ soit ample sur $X_V/V$.
\end{enumerate}
Il est clair que (A) entraîne le lemme.

\medskip\noindent
Soient $X \to Y$, $\mathcal{L}$ et $\Spec(k) \to Y$ comme dans (A).
D'après le lemme \ref{lemma-descending-property-ample},
on peut supposer que $k = \kappa(y)$ est le corps résiduel d'un point $y$
de $Y$.

\medskip\noindent
Comme $Y$ est affine, on peut trouver un ensemble préordonné filtrant $I$ et un
système projectif de morphismes $X_i \to Y_i$ d'espaces algébriques,
où $Y_i$ est de présentation finie sur $\mathbf{Z}$, où les
morphismes de transition $X_i \to X_{i'}$ et $Y_i \to Y_{i'}$ sont affines,
où $X_i \to Y_i$ est propre et de présentation finie, et tel que
$X \to Y = \lim (X_i \to Y_i)$. Voir Limites d'espaces, lemme
\ref{spaces-limits-lemma-proper-limit-of-proper-finite-presentation-noetherian}.
Quitte à remplacer $I$ par une partie cofinale, on peut supposer que $Y_i$ est un schéma (affine) pour tout $i$ ;
voir Limites d'espaces, lemme \ref{spaces-limits-lemma-limit-is-affine}.
Quitte encore à remplacer $I$ par une partie cofinale, on peut supposer disposer d'un système compatible de
$\mathcal{O}_{X_i}$-modules inversibles $\mathcal{L}_i$
dont l'image inverse est $\mathcal{L}$ ; voir
Limites d'espaces, lemme \ref{spaces-limits-lemma-descend-invertible-modules}.
Soit $y_i \in Y_i$ l'image de $y$.
Alors $\kappa(y) = \colim \kappa(y_i)$.
Par conséquent, $X_y = \lim X_{i, y_i}$ et, quitte à remplacer $I$ par une partie cofinale,
on peut supposer que $X_{i, y_i}$ est un schéma pour tout $i$ ; voir
Limites d'espaces, lemme \ref{spaces-limits-lemma-limit-is-scheme}.
Il existe donc un $i$ tel que $\mathcal{L}_{i, y_i}$
soit ample sur $X_{i, y_i}$, d'après
Limites, lemme \ref{limits-lemma-limit-ample}.
D'après Diviseurs sur les espaces, lemme \ref{spaces-divisors-lemma-ample-in-neighbourhood},
on trouve un voisinage ouvert
$V_i \subset Y_i$ de $y_i$ tel que
la restriction de $\mathcal{L}_i$ à $f_i^{-1}(V_i)$
soit ample relativement à $V_i$.
On achève la démonstration en prenant pour $V \subset Y$ l'image inverse de
$V_i$ (indications : utiliser
Morphismes, lemme \ref{morphisms-lemma-ample-base-change},
le fait que $X \to Y \times_{Y_i} X_i$ est affine,
et le fait que l'image inverse d'un
faisceau inversible ample par un morphisme affine est ample, d'après
Morphismes, lemme \ref{morphisms-lemma-pullback-ample-tensor-relatively-ample}).
\end{proof}













\section{Propriétés des morphismes locales sur la source}
\label{section-properties-morphisms-local-source}

\noindent
Dans cette section, nous définissons ce que signifie, pour une propriété des morphismes
d'espaces algébriques, être locale sur la source. Comparer avec
Descente, section \ref{descent-section-properties-morphisms-local-source}.

\begin{definition}
\label{definition-property-morphisms-local-source}
Soit $S$ un schéma.
Soit $\mathcal{P}$ une propriété des morphismes d'espaces algébriques sur $S$.
Soit $\tau \in \{fpqc, \linebreak[0] fppf, \linebreak[0] syntomic, \linebreak[0]
smooth, \linebreak[0] \etale\}$. Nous disons que $\mathcal{P}$ est
{\it $\tau$-locale sur la source}, ou
{\it locale sur la source pour la topologie $\tau$}, si, pour
tout morphisme $f : X \to Y$ d'espaces algébriques sur $S$ et tout
$\tau$-recouvrement $\{X_i \to X\}_{i \in I}$ d'espaces algébriques, on a
$$
f \text{ vérifie }\mathcal{P}
\Leftrightarrow
\text{chaque }X_i \to Y\text{ vérifie }\mathcal{P}.
$$
\end{definition}

\noindent
Précisons que, puisque les isomorphismes sont toujours couvrants,
on voit (ou l'on exige) que la propriété $\mathcal{P}$ est vérifiée par $X \to Y$
si et seulement si elle est vérifiée par toute flèche $X' \to Y'$ isomorphe à $X \to Y$.
Si une propriété est $\tau$-locale sur la source, elle reste vérifiée après
précomposition par tout morphisme figurant dans un $\tau$-recouvrement. Voici
un énoncé formel.

\begin{lemma}
\label{lemma-precompose-property-local-source}
Soit $S$ un schéma.
Soit $\tau \in \{fpqc, \linebreak[0] fppf, \linebreak[0] syntomic, \linebreak[0]
smooth, \linebreak[0] \etale\}$.
Soit $\mathcal{P}$ une propriété des morphismes d'espaces algébriques sur $S$
qui est $\tau$-locale sur la source. Supposons que $f : X \to Y$ vérifie la propriété
$\mathcal{P}$. Pour tout morphisme $a : X' \to X$ qui est respectivement
plat, plat et localement de présentation finie, syntomique,
lisse ou \'etale, la composée $f \circ a : X' \to Y$ vérifie
la propriété $\mathcal{P}$.
\end{lemma}

\begin{proof}
Cela résulte du fait que l'on peut compléter $X' \to X$ en une famille de
morphismes formant un recouvrement pour la topologie $\tau$.
\end{proof}

\begin{lemma}
\label{lemma-transfer-from-schemes}
Soit $S$ un schéma.
Soit $\tau \in \{fpqc, \linebreak[0] fppf, \linebreak[0] syntomic, \linebreak[0]
smooth, \linebreak[0] \etale\}$.
Supposons que $\mathcal{P}$ soit une propriété des morphismes de schémas sur $S$
qui soit locale sur la source et le but pour la topologie \'etale. Notons $\mathcal{P}_{spaces}$
la propriété correspondante des morphismes d'espaces algébriques sur $S$ ; voir
Morphismes d'espaces, définition \ref{spaces-morphisms-definition-P}.
Si $\mathcal{P}$ est locale sur la source pour la topologie $\tau$, alors
$\mathcal{P}_{spaces}$ est locale sur la source pour la topologie $\tau$.
\end{lemma}

\begin{proof}
Soit $f : X \to Y$ un morphisme d'espaces algébriques sur $S$.
Soit $\{X_i \to X\}_{i \in I}$ un recouvrement d'espaces algébriques pour la topologie $\tau$.
Choisissons un schéma $V$ et un morphisme \'etale surjectif $V \to Y$.
Choisissons un schéma $U$ et un morphisme \'etale surjectif $U \to X \times_Y V$.
Pour chaque $i$, choisissons un schéma $U_i$ et un morphisme \'etale surjectif
$U_i \to X_i \times_X U$.

\medskip\noindent
Remarquons que $\{X_i \times_X U \to U\}_{i \in I}$ est un recouvrement pour la topologie $\tau$.
Remarquons que chacune des familles $\{U_i \to X_i \times_X U\}$ est un recouvrement \'etale,
donc un recouvrement pour la topologie $\tau$. Ainsi, $\{U_i \to U\}_{i \in I}$ est un
recouvrement pour la topologie $\tau$ dans la catégorie des espaces algébriques sur $S$. Mais puisque $U$ et chaque $U_i$
sont des schémas, $\{U_i \to U\}_{i \in I}$ est un
recouvrement pour la topologie $\tau$ dans la catégorie des schémas sur $S$.

\medskip\noindent
On a maintenant
\begin{align*}
f \text{ vérifie }\mathcal{P}_{spaces}
& \Leftrightarrow
U \to V \text{ vérifie }\mathcal{P} \\
& \Leftrightarrow
\text{chaque }U_i \to V \text{ vérifie }\mathcal{P} \\
& \Leftrightarrow
\text{chaque }X_i \to Y\text{ vérifie }\mathcal{P}_{spaces}.
\end{align*}
La première et la dernière équivalence résultent de la définition de
$\mathcal{P}_{spaces}$, tandis que l'équivalence intermédiaire résulte de l'hypothèse
selon laquelle $\mathcal{P}$ est locale sur la source pour la topologie $\tau$.
\end{proof}
% preserved blank authority line 2405
% preserved blank authority line 2406
% preserved blank authority line 2407
% preserved blank authority line 2408
% preserved blank authority line 2409
% preserved blank authority line 2410
\section{Propriétés des morphismes locales sur la source pour la topologie fpqc}
\label{section-fpqc-local-source}

\noindent
Voici quelques propriétés des morphismes qui sont locales sur la source pour la topologie fpqc.

\begin{lemma}
\label{lemma-flat-fpqc-local-source}
La propriété $\mathcal{P}(f)=$``$f$ est plat'' est locale sur la source pour la topologie fpqc.
\end{lemma}

\begin{proof}
L'assertion résulte du
lemme \ref{lemma-transfer-from-schemes}
en utilisant
Morphismes d'espaces, définition \ref{spaces-morphisms-definition-flat}
et
Descente, lemme \ref{descent-lemma-flat-fpqc-local-source}.
\end{proof}










\section{Propriétés des morphismes locales sur la source pour la topologie fppf}
\label{section-fppf-local-source}

\noindent
Voici quelques propriétés des morphismes qui sont locales sur la source pour la topologie fppf.

\begin{lemma}
\label{lemma-locally-finite-presentation-fppf-local-source}
La propriété $\mathcal{P}(f)=$``$f$ est localement de présentation finie''
est locale sur la source pour la topologie fppf.
\end{lemma}

\begin{proof}
L'assertion résulte du
lemme \ref{lemma-transfer-from-schemes}
en utilisant
Morphismes d'espaces,
définition \ref{spaces-morphisms-definition-locally-finite-presentation}
et
Descente,
lemme \ref{descent-lemma-locally-finite-presentation-fppf-local-source}.
\end{proof}

\begin{lemma}
\label{lemma-locally-finite-type-fppf-local-source}
La propriété $\mathcal{P}(f)=$``$f$ est localement de type fini''
est locale sur la source pour la topologie fppf.
\end{lemma}

\begin{proof}
L'assertion résulte du
lemme \ref{lemma-transfer-from-schemes}
en utilisant
Morphismes d'espaces,
définition \ref{spaces-morphisms-definition-locally-finite-type}
et
Descente, lemme \ref{descent-lemma-locally-finite-type-fppf-local-source}.
\end{proof}

\begin{lemma}
\label{lemma-open-fppf-local-source}
La propriété $\mathcal{P}(f)=$``$f$ est ouvert''
est locale sur la source pour la topologie fppf.
\end{lemma}

\begin{proof}
L'assertion résulte du
lemme \ref{lemma-transfer-from-schemes}
en utilisant
Morphismes d'espaces, définition \ref{spaces-morphisms-definition-open}
et
Descente, lemme \ref{descent-lemma-open-fppf-local-source}.
\end{proof}

\begin{lemma}
\label{lemma-universally-open-fppf-local-source}
La propriété $\mathcal{P}(f)=$``$f$ est universellement ouvert''
est locale sur la source pour la topologie fppf.
\end{lemma}

\begin{proof}
L'assertion résulte du
lemme \ref{lemma-transfer-from-schemes}
en utilisant
Morphismes d'espaces, définition \ref{spaces-morphisms-definition-open}
et
Descente, lemme \ref{descent-lemma-universally-open-fppf-local-source}.
\end{proof}



\section{Propriétés des morphismes locales sur la source pour la topologie syntomique}
\label{section-syntomic-local-source}

\noindent
Voici quelques propriétés des morphismes qui sont locales sur la source pour la topologie syntomique.

\begin{lemma}
\label{lemma-syntomic-syntomic-local-source}
La propriété $\mathcal{P}(f)=$``$f$ est syntomique''
est locale sur la source pour la topologie syntomique.
\end{lemma}

\begin{proof}
L'assertion résulte du
lemme \ref{lemma-transfer-from-schemes}
en utilisant
Morphismes d'espaces, définition \ref{spaces-morphisms-definition-syntomic}
et
Descente, lemme \ref{descent-lemma-syntomic-syntomic-local-source}.
\end{proof}




\section{Propriétés des morphismes locales sur la source pour la topologie lisse}
\label{section-smooth-local-source}

\noindent
Voici quelques propriétés des morphismes qui sont locales sur la source pour la topologie lisse.

\begin{lemma}
\label{lemma-smooth-smooth-local-source}
La propriété $\mathcal{P}(f)=$``$f$ est lisse''
est locale sur la source pour la topologie lisse.
\end{lemma}

\begin{proof}
L'assertion résulte du
lemme \ref{lemma-transfer-from-schemes}
en utilisant
Morphismes d'espaces, définition \ref{spaces-morphisms-definition-smooth}
et
Descente, lemme \ref{descent-lemma-smooth-smooth-local-source}.
\end{proof}



\section{Propriétés des morphismes locales sur la source pour la topologie \'etale}
\label{section-etale-local-source}

\noindent
Voici quelques propriétés des morphismes qui sont locales sur la source pour la topologie \'etale.

\begin{lemma}
\label{lemma-etale-etale-local-source}
La propriété $\mathcal{P}(f)=$``$f$ est \'etale''
est locale sur la source pour la topologie \'etale.
\end{lemma}

\begin{proof}
L'assertion résulte du
lemme \ref{lemma-transfer-from-schemes}
en utilisant
Morphismes d'espaces,
définition \ref{spaces-morphisms-definition-etale}
et
Descente, lemme \ref{descent-lemma-etale-etale-local-source}.
\end{proof}

\begin{lemma}
\label{lemma-locally-quasi-finite-etale-local-source}
La propriété $\mathcal{P}(f)=$``$f$ est localement quasi-fini''
est locale sur la source pour la topologie \'etale.
\end{lemma}

\begin{proof}
L'assertion résulte du
lemme \ref{lemma-transfer-from-schemes}
en utilisant
Morphismes d'espaces,
définition \ref{spaces-morphisms-definition-locally-quasi-finite}
et
Descente, lemme \ref{descent-lemma-locally-quasi-finite-etale-local-source}.
\end{proof}

\begin{lemma}
\label{lemma-unramified-etale-local-source}
La propriété $\mathcal{P}(f)=$``$f$ est non ramifié''
est locale sur la source pour la topologie \'etale.
\end{lemma}

\begin{proof}
L'assertion résulte du
lemme \ref{lemma-transfer-from-schemes}
en utilisant
Morphismes d'espaces, définition \ref{spaces-morphisms-definition-unramified}
et
Descente, lemme \ref{descent-lemma-unramified-etale-local-source}.
\end{proof}


% preserved blank authority line 2612
% preserved blank authority line 2613
\section{Propriétés des morphismes locales pour la topologie lisse sur la source et le but}
\sectionmark{Propriétés locales pour la topologie lisse}
\label{section-properties-local-source-target}

\noindent
Soit $\mathcal{P}$ une propriété des morphismes d'espaces algébriques. La phrase
« $\mathcal{P}$ est locale pour la topologie lisse sur la
source et le but » possède un sens intuitif. Toutefois, cette notion ne revient
pas à demander que $\mathcal{P}$ soit à la fois locale pour la topologie lisse
sur la source et locale pour la topologie lisse sur le but.
Nous avons étudié très en détail un phénomène analogue (pour la topologie \'etale
et la catégorie des schémas) dans
Descente, section \ref{descent-section-properties-etale-local-source-target}
(pour un aperçu rapide, voir
Descente, remarque \ref{descent-remark-compare-definitions}).
Il existe toutefois une différence importante entre le cas de la topologie lisse
et celui de la topologie \'etale. Pour la saisir, nous invitons le lecteur
à comparer
Descente, lemme \ref{descent-lemma-local-source-target-implies}
et
le lemme \ref{lemma-local-source-target-implies},
ainsi qu'à comparer
Descente, lemme \ref{descent-lemma-local-source-target-characterize}
et
le lemme \ref{lemma-local-source-target-characterize}.
En effet, dans le cadre étale, le choix du « recouvrement » étale du
but est indifférent, tandis qu'il ne l'est pas dans le cadre lisse.

\begin{definition}
\label{definition-local-source-target}
Soit $S$ un schéma.
Soit $\mathcal{P}$ une propriété des morphismes d'espaces algébriques sur $S$.
Nous disons que $\mathcal{P}$ est {\it locale pour la topologie lisse sur la source et le but} si
\begin{enumerate}
\item (stabilité par précomposition par des morphismes lisses)
si $f : X \to Y$ est lisse et si $g : Y \to Z$ vérifie $\mathcal{P}$,
alors $g \circ f$ vérifie $\mathcal{P}$,
\item (stabilité par changement de base lisse)
si $f : X \to Y$ vérifie $\mathcal{P}$ et si $Y' \to Y$ est lisse, alors
le changement de base $f' : Y' \times_Y X \to Y'$ vérifie $\mathcal{P}$, et
\item (caractère local) pour un morphisme $f : X \to Y$, les conditions suivantes sont
équivalentes :
\begin{enumerate}
\item $f$ vérifie $\mathcal{P}$,
\item pour tout $x \in |X|$, il existe un diagramme commutatif
$$
\xymatrix{
U \ar[d]_a \ar[r]_h & V \ar[d]^b \\
X \ar[r]^f & Y
}
$$
dont les flèches verticales sont lisses et pour lequel il existe $u \in |U|$ tel que $a(u) = x$ et que
$h$ vérifie $\mathcal{P}$.
\end{enumerate}
\end{enumerate}
\end{definition}

\noindent
Ce qui précède nous sert de définition. Dans les lemmes ci-dessous, nous montrerons que
cela équivaut à ce que $\mathcal{P}$ soit locale pour la topologie lisse sur le but,
locale pour la topologie lisse sur la source et stable par postcomposition par des morphismes lisses.

\begin{lemma}
\label{lemma-local-source-target-implies}
Soit $S$ un schéma.
Soit $\mathcal{P}$ une propriété des morphismes d'espaces algébriques sur $S$
qui est locale pour la topologie lisse sur la source et le but. Alors :
\begin{enumerate}
\item $\mathcal{P}$ est locale pour la topologie lisse sur la source,
\item $\mathcal{P}$ est locale pour la topologie lisse sur le but,
\item $\mathcal{P}$ est stable par postcomposition par des morphismes lisses :
si $f : X \to Y$ vérifie $\mathcal{P}$ et si $g : Y \to Z$ est lisse, alors
$g \circ f$ vérifie $\mathcal{P}$.
\end{enumerate}
\end{lemma}

\begin{proof}
Écrivons tous les détails.

\medskip\noindent
Preuve de (1). Soit $f : X \to Y$ un morphisme d'espaces algébriques sur $S$.
Soit $\{X_i \to X\}_{i \in I}$ un recouvrement lisse de $X$. Si chaque composé
$h_i : X_i \to Y$ vérifie $\mathcal{P}$, alors, pour tout $|x| \in X$, on peut trouver
un indice $i \in I$ et un point $x_i \in |X_i|$ qui s'envoie sur $x$. Alors
$(X_i, x_i) \to (X, x)$ est un morphisme lisse de couples,
$\text{id}_Y : Y \to Y$ est un morphisme lisse et $h_i$ est comme dans la partie (3) de la
définition \ref{definition-local-source-target}.
On voit donc que $f$ vérifie $\mathcal{P}$.
Réciproquement, si $f$ vérifie $\mathcal{P}$, alors chaque $X_i \to Y$ vérifie
$\mathcal{P}$ d'après la
partie (1) de la définition \ref{definition-local-source-target}.

\medskip\noindent
Preuve de (2). Soit $f : X \to Y$ un morphisme d'espaces algébriques sur $S$.
Soit $\{Y_i \to Y\}_{i \in I}$ un recouvrement lisse de $Y$.
Notons $X_i = Y_i \times_Y X$ et $h_i : X_i \to Y_i$ le changement de base
de $f$. Si chaque $h_i : X_i \to Y_i$ vérifie $\mathcal{P}$, alors, pour tout
$x \in |X|$, choisissons un indice $i \in I$ et un point $x_i \in |X_i|$ qui s'envoie sur $x$.
Alors $(X_i, x_i) \to (X, x)$ est un morphisme lisse de couples, $Y_i \to Y$ est
lisse et $h_i$ est comme dans la partie (3) de la
définition \ref{definition-local-source-target}.
On voit donc que $f$ vérifie $\mathcal{P}$.
Réciproquement, si $f$ vérifie $\mathcal{P}$, alors chaque $X_i \to Y_i$ vérifie
$\mathcal{P}$ d'après la
partie (2) de la définition \ref{definition-local-source-target}.

\medskip\noindent
Preuve de (3). Supposons que $f : X \to Y$ vérifie $\mathcal{P}$ et que $g : Y \to Z$ soit
lisse. Pour tout $x \in |X|$, considérons $(X, x) \to (X, x)$ comme un
morphisme lisse de couples ; $Y \to Z$ est un morphisme lisse, et $h = f$ est comme
dans la partie (3) de la
définition \ref{definition-local-source-target}.
On voit donc que $g \circ f$ vérifie $\mathcal{P}$.
\end{proof}

\noindent
Le lemme suivant est l'analogue de
Morphismes, lemme \ref{morphisms-lemma-locally-P-characterize}.

\begin{lemma}
\label{lemma-local-source-target-characterize}
Soit $S$ un schéma.
Soit $\mathcal{P}$ une propriété des morphismes d'espaces algébriques sur $S$
qui est locale pour la topologie lisse sur la source et le but. Soit $f : X \to Y$ un morphisme
d'espaces algébriques sur $S$. Les conditions suivantes sont équivalentes :
\begin{enumerate}
\item[(a)] $f$ vérifie la propriété $\mathcal{P}$,
\item[(b)] pour tout $x \in |X|$, il existe un morphisme lisse de couples
$a : (U, u) \to (X, x)$, un morphisme lisse $b : V \to Y$ et
un morphisme $h : U \to V$ tel que $f \circ a = b \circ h$ et que
$h$ vérifie $\mathcal{P}$,
\item[(c)] pour un certain diagramme commutatif
$$
\xymatrix{
U \ar[d]_a \ar[r]_h & V \ar[d]^b \\
X \ar[r]^f & Y
}
$$
où $a$ et $b$ sont lisses et $a$ est surjectif, le morphisme $h$ vérifie $\mathcal{P}$,
\item[(d)] pour tout diagramme commutatif
$$
\xymatrix{
U \ar[d]_a \ar[r]_h & V \ar[d]^b \\
X \ar[r]^f & Y
}
$$
où $b$ est lisse et $U \to X \times_Y V$ est lisse,
le morphisme $h$ vérifie $\mathcal{P}$,
\item[(e)] il existe un recouvrement lisse $\{Y_i \to Y\}_{i \in I}$ tel
que chaque changement de base $Y_i \times_Y X \to Y_i$ vérifie $\mathcal{P}$,
\item[(f)] il existe un recouvrement lisse $\{X_i \to X\}_{i \in I}$ tel
que chaque composé $X_i \to Y$ vérifie $\mathcal{P}$,
\item[(g)] il existe un recouvrement lisse $\{Y_i \to Y\}_{i \in I}$ et,
pour chaque $i \in I$, un recouvrement lisse
$\{X_{ij} \to Y_i \times_Y X\}_{j \in J_i}$ tel que chaque morphisme
$X_{ij} \to Y_i$ vérifie $\mathcal{P}$.
\end{enumerate}
\end{lemma}

\begin{proof}
L'équivalence de (a) et (b) fait partie de la
définition \ref{definition-local-source-target}.
L'équivalence de (a) et (e) est donnée par la
partie (2) du lemme \ref{lemma-local-source-target-implies}.
L'équivalence de (a) et (f) est donnée par la
partie (1) du lemme \ref{lemma-local-source-target-implies}.
Puisque (a) est maintenant équivalente à (e) et à (f), il s'ensuit que
(a) est équivalente à (g).

\medskip\noindent
Il est clair que (c) implique (b). Si (b) est vérifiée, alors, pour tout
$x \in |X|$, nous pouvons choisir un morphisme lisse de couples
$a_x : (U_x, u_x) \to (X, x)$, un morphisme lisse $b_x : V_x \to Y$ et
un morphisme $h_x : U_x \to V_x$ tel que $f \circ a_x = b_x \circ h_x$ et que
$h_x$ vérifie $\mathcal{P}$. Alors $h = \coprod h_x : \coprod U_x \to \coprod V_x$
forme, avec $a = \coprod a_x$ et $b = \coprod b_x$, un diagramme comme en (c).
(Remarquons que $h$ vérifie la propriété $\mathcal{P}$, puisque $\{V_x \to \coprod V_x\}$
est un recouvrement lisse et que $\mathcal{P}$ est locale pour la topologie lisse sur le but.)
Ainsi, (b) et (c) sont équivalentes.

\medskip\noindent
Nous savons maintenant que (a), (b), (c), (e), (f) et (g) sont équivalentes.
Supposons (a) vérifiée. Soient $U, V, a, b, h$ comme en (d). Alors
$X \times_Y V \to V$ vérifie $\mathcal{P}$ puisque $\mathcal{P}$ est stable par
changement de base lisse ; il s'ensuit que $U \to V$ vérifie $\mathcal{P}$ puisque $\mathcal{P}$
est stable par précomposition par des morphismes lisses. Réciproquement, si (d)
est vérifiée, prendre $U = X$ et $V = Y$ montre que $f$ vérifie $\mathcal{P}$.
\end{proof}

\begin{lemma}
\label{lemma-smooth-local-source-target}
Soit $S$ un schéma.
Soit $\mathcal{P}$ une propriété des morphismes d'espaces algébriques sur $S$.
Supposons que
\begin{enumerate}
\item $\mathcal{P}$ soit locale pour la topologie lisse sur la source,
\item $\mathcal{P}$ soit locale pour la topologie lisse sur le but, et
\item $\mathcal{P}$ soit stable par postcomposition par des morphismes lisses :
si $f : X \to Y$ vérifie $\mathcal{P}$ et si $Y \to Z$ est un morphisme lisse,
alors $X \to Z$ vérifie $\mathcal{P}$.
\end{enumerate}
Alors $\mathcal{P}$ est locale pour la topologie lisse sur la source et le but.
\end{lemma}

\begin{proof}
Soit $\mathcal{P}$ une propriété des morphismes d'espaces algébriques qui
satisfait aux conditions (1), (2) et (3) du lemme. D'après le
lemme \ref{lemma-precompose-property-local-source},
on voit que $\mathcal{P}$ est stable par précomposition par des
morphismes lisses. D'après le
lemme \ref{lemma-pullback-property-local-target},
on voit que $\mathcal{P}$ est stable par changement de base lisse.
Il suffit donc de prouver que la partie (3) de la
définition \ref{definition-local-source-target}
est vérifiée.

\medskip\noindent
Plus précisément, supposons que $f : X \to Y$ soit un morphisme
d'espaces algébriques sur $S$ qui vérifie la
partie (3)(b) de la définition \ref{definition-local-source-target}.
Autrement dit, pour tout $x \in X$, il existe un morphisme lisse
$a_x : U_x \to X$, un point $u_x \in |U_x|$ qui s'envoie sur $x$,
un morphisme lisse $b_x : V_x \to Y$ et un morphisme $h_x : U_x \to V_x$
tels que $f \circ a_x = b_x \circ h_x$ et que $h_x$ vérifie $\mathcal{P}$.
Pour achever la preuve du lemme, il suffit de montrer que $f$ vérifie $\mathcal{P}$.
Posons $U = \coprod U_x$, $a = \coprod a_x$, $V = \coprod V_x$,
$b = \coprod b_x$ et $h = \coprod h_x$. Nous obtenons un
diagramme commutatif
$$
\xymatrix{
U \ar[d]_a \ar[r]_h & V \ar[d]^b \\
X \ar[r]^f & Y
}
$$
où $a$ et $b$ sont lisses et $a$ est surjectif. Remarquons que $h$ vérifie $\mathcal{P}$
puisque chaque $h_x$ la vérifie et que $\mathcal{P}$ est locale pour la topologie lisse sur le but.
Comme $a$ est surjectif et que $\mathcal{P}$ est locale pour la topologie lisse sur la source,
il suffit de prouver que $b \circ h$ vérifie $\mathcal{P}$.
Cela résulte de l'hypothèse selon laquelle $\mathcal{P}$ est stable par
postcomposition par un morphisme lisse et du fait que $b$ est lisse.
\end{proof}

\begin{remark}
\label{remark-list-local-source-target}
À l'aide du
lemme \ref{lemma-smooth-local-source-target}
et des résultats établis dans les sections précédentes de ce chapitre, on dresse aisément
une liste de types de morphismes dont la propriété est locale pour la topologie lisse sur la
source et le but. Dans chaque cas, nous indiquons le lemme qui implique que
la propriété est locale pour la topologie lisse sur la source et celui qui implique qu'elle
est locale pour la topologie lisse sur le but. Dans chaque cas, la troisième hypothèse
du
lemme \ref{lemma-smooth-local-source-target}
se vérifie immédiatement ; nous l'omettons. Voici la liste :
\begin{enumerate}
\item plat, voir les
lemmes \ref{lemma-flat-fpqc-local-source} et
\ref{lemma-descending-property-flat},
\item localement de présentation finie, voir les
lemmes \ref{lemma-locally-finite-presentation-fppf-local-source} et
\ref{lemma-descending-property-locally-finite-presentation},
\item localement de type fini, voir les
lemmes \ref{lemma-locally-finite-type-fppf-local-source} et
\ref{lemma-descending-property-locally-finite-type},
\item universellement ouvert, voir les
lemmes \ref{lemma-universally-open-fppf-local-source} et
\ref{lemma-descending-property-universally-open},
\item syntomique, voir les
lemmes \ref{lemma-syntomic-syntomic-local-source} et
\ref{lemma-descending-property-syntomic},
\item lisse, voir les
lemmes \ref{lemma-smooth-smooth-local-source} et
\ref{lemma-descending-property-smooth},
\item ajouter ici d'autres cas au besoin.
\end{enumerate}
\end{remark}






% CH074 slice boundary 2895
\section{Propriétés des morphismes locales pour la topologie \'etale sur la source et la topologie lisse sur le but}
\sectionmark{Topologie \'etale sur la source, lisse sur le but}
\label{section-properties-etale-smooth-local-source-target}

\noindent
Cette section est l'analogue de la
section \ref{section-properties-local-source-target}
pour les propriétés des morphismes qui sont locales pour la topologie \'etale
sur la source et pour la topologie lisse sur le but.
Nous donnons à cette propriété un nom ridiculement long
afin de ne pas avoir à l'employer trop souvent.

\begin{definition}
\label{definition-etale-smooth-local-source-target}
Soit $S$ un schéma.
Soit $\mathcal{P}$ une propriété des morphismes d'espaces algébriques sur $S$.
Nous disons que $\mathcal{P}$ est {\it locale pour la topologie \'etale sur la source et pour la topologie lisse sur le but} si
\begin{enumerate}
\item (stabilité par précomposition par des morphismes \'etales)
si $f : X \to Y$ est \'etale et si $g : Y \to Z$ vérifie $\mathcal{P}$,
alors $g \circ f$ vérifie $\mathcal{P}$,
\item (stabilité par changement de base lisse)
si $f : X \to Y$ vérifie $\mathcal{P}$ et si $Y' \to Y$ est lisse, alors
le changement de base $f' : Y' \times_Y X \to Y'$ vérifie $\mathcal{P}$, et
\item (caractère local) pour un morphisme $f : X \to Y$, les conditions suivantes sont
équivalentes :
\begin{enumerate}
\item $f$ vérifie $\mathcal{P}$,
\item pour tout $x \in |X|$, il existe un diagramme commutatif
$$
\xymatrix{
U \ar[d]_a \ar[r]_h & V \ar[d]^b \\
X \ar[r]^f & Y
}
$$
où $b$ est lisse, où $U \to X \times_Y V$ est \'etale
et où il existe $u \in |U|$ tel que $a(u) = x$ et que
$h$ vérifie $\mathcal{P}$.
\end{enumerate}
\end{enumerate}
\end{definition}

\noindent
Ce qui précède nous sert de définition. Dans les lemmes ci-dessous, nous montrerons que
cela équivaut à ce que $\mathcal{P}$ soit locale pour la topologie \'etale sur le but,
locale pour la topologie lisse sur la source et stable par postcomposition par des
morphismes \'etales.

\begin{lemma}
\label{lemma-etale-smooth-local-source-target-implies}
Soit $S$ un schéma.
Soit $\mathcal{P}$ une propriété des morphismes d'espaces algébriques sur $S$
qui est locale pour la topologie \'etale sur la source et pour la topologie lisse sur le but. Alors :
\begin{enumerate}
\item $\mathcal{P}$ est locale pour la topologie \'etale sur la source,
\item $\mathcal{P}$ est locale pour la topologie lisse sur le but,
\item $\mathcal{P}$ est stable par postcomposition par des morphismes \'etales :
si $f : X \to Y$ vérifie $\mathcal{P}$ et si $g : Y \to Z$ est \'etale, alors
$g \circ f$ vérifie $\mathcal{P}$, et
\item $\mathcal{P}$ possède la propriété de permanence suivante : si $f : X \to Y$ et
$g : Y \to Z$ est \'etale et si $g \circ f$ vérifie $\mathcal{P}$, alors
$f$ vérifie $\mathcal{P}$.
\end{enumerate}
\end{lemma}

\begin{proof}
Écrivons tous les détails.

\medskip\noindent
Preuve de (1). Soit $f : X \to Y$ un morphisme d'espaces algébriques sur $S$.
Soit $\{X_i \to X\}_{i \in I}$ un recouvrement \'etale de $X$. Si chaque composé
$h_i : X_i \to Y$ vérifie $\mathcal{P}$, alors, pour tout $|x| \in X$, on peut trouver
un indice $i \in I$ et un point $x_i \in |X_i|$ qui s'envoie sur $x$. Alors
$(X_i, x_i) \to (X, x)$ est un morphisme \'etale de couples,
$\text{id}_Y : Y \to Y$ est un morphisme lisse et $h_i$ est comme dans la partie (3) de la
définition \ref{definition-etale-smooth-local-source-target}.
On voit donc que $f$ vérifie $\mathcal{P}$.
Réciproquement, si $f$ vérifie $\mathcal{P}$, alors chaque $X_i \to Y$ vérifie
$\mathcal{P}$ d'après la
partie (1) de la définition \ref{definition-etale-smooth-local-source-target}.

\medskip\noindent
Preuve de (2). Soit $f : X \to Y$ un morphisme d'espaces algébriques sur $S$.
Soit $\{Y_i \to Y\}_{i \in I}$ un recouvrement lisse de $Y$.
Notons $X_i = Y_i \times_Y X$ et $h_i : X_i \to Y_i$ le changement de base
de $f$. Si chaque $h_i : X_i \to Y_i$ vérifie $\mathcal{P}$, alors, pour tout
$x \in |X|$, choisissons un indice $i \in I$ et un point $x_i \in |X_i|$ qui s'envoie sur $x$.
Alors $X_i \to X \times_Y Y_i$ est un morphisme \'etale
(puisque c'est un isomorphisme), $Y_i \to Y$ est
lisse et $h_i$ est comme dans la partie (3) de la
définition \ref{definition-local-source-target}.
On voit donc que $f$ vérifie $\mathcal{P}$.
Réciproquement, si $f$ vérifie $\mathcal{P}$, alors chaque $X_i \to Y_i$ vérifie
$\mathcal{P}$ d'après la
partie (2) de la définition \ref{definition-local-source-target}.

\medskip\noindent
Preuve de (3). Supposons que $f : X \to Y$ vérifie $\mathcal{P}$ et que $g : Y \to Z$ soit
\'etale. Le morphisme $X \to Y \times_Z X$ est \'etale en tant que morphisme
entre des espaces algébriques \'etales sur $X$ (
Propriétés des espaces, lemme \ref{spaces-properties-lemma-etale-permanence}).
De plus, $Y \to Z$ est \'etale, donc lisse.
Ainsi, le diagramme
$$
\xymatrix{
X \ar[d] \ar[r]_f & Y \ar[d] \\
X \ar[r]^{g \circ f} & Z
}
$$
convient, pour tout $x \in |X|$, dans la partie (3) de la
définition \ref{definition-local-source-target},
et nous concluons que $g \circ f$ vérifie $\mathcal{P}$.

\medskip\noindent
Preuve de (4). Soient $f : X \to Y$ un morphisme et $g : Y \to Z$ \'etale
tels que $g \circ f$ vérifie $\mathcal{P}$. Alors, d'après la
partie (2) de la définition \ref{definition-etale-smooth-local-source-target},
on voit que $\text{pr}_Y : Y \times_Z X \to Y$ vérifie $\mathcal{P}$. Mais
le morphisme $(f, 1) : X \to Y \times_Z X$ est \'etale, car c'est une section de la
projection \'etale $\text{pr}_X : Y \times_Z X \to X$ ; voir
Morphismes d'espaces, lemme \ref{spaces-morphisms-lemma-etale-permanence}.
Ainsi, $f = \text{pr}_Y \circ (f, 1)$ vérifie $\mathcal{P}$ d'après la
partie (1) de la définition \ref{definition-etale-smooth-local-source-target}.
\end{proof}

\begin{lemma}
\label{lemma-etale-smooth-local-source-target-characterize}
Soit $S$ un schéma.
Soit $\mathcal{P}$ une propriété des morphismes d'espaces algébriques sur $S$
qui est locale pour la topologie étale sur la source et pour la topologie lisse sur le but.
Soit $f : X \to Y$ un morphisme
d'espaces algébriques sur $S$. Les conditions suivantes sont équivalentes :
\begin{enumerate}
\item[(a)] $f$ vérifie la propriété $\mathcal{P}$,
\item[(b)] pour tout $x \in |X|$, il existe un morphisme lisse $b : V \to Y$,
un morphisme \'etale $a : U \to V \times_Y X$ et un point $u \in |U|$
qui s'envoie sur $x$, de telle sorte que $U \to V$ vérifie $\mathcal{P}$,
\item[(c)] pour un certain diagramme commutatif
$$
\xymatrix{
U \ar[d]_a \ar[r]_h & V \ar[d]^b \\
X \ar[r]^f & Y
}
$$
où $b$ est lisse, $U \to V \times_Y X$ est \'etale et $a$ est surjectif,
le morphisme $h$ vérifie $\mathcal{P}$,
\item[(d)] pour tout diagramme commutatif
$$
\xymatrix{
U \ar[d]_a \ar[r]_h & V \ar[d]^b \\
X \ar[r]^f & Y
}
$$
où $b$ est lisse et $U \to X \times_Y V$ est \'etale, le morphisme $h$
vérifie $\mathcal{P}$,
\item[(e)] il existe un recouvrement lisse $\{Y_i \to Y\}_{i \in I}$ tel
que chaque changement de base $Y_i \times_Y X \to Y_i$ vérifie $\mathcal{P}$,
\item[(f)] il existe un recouvrement \'etale $\{X_i \to X\}_{i \in I}$ tel
que chaque composé $X_i \to Y$ vérifie $\mathcal{P}$,
\item[(g)] il existe un recouvrement lisse $\{Y_i \to Y\}_{i \in I}$ et,
pour chaque $i \in I$, un recouvrement \'etale
$\{X_{ij} \to Y_i \times_Y X\}_{j \in J_i}$ tel que chaque morphisme
$X_{ij} \to Y_i$ vérifie $\mathcal{P}$.
\end{enumerate}
\end{lemma}

\begin{proof}
L'équivalence de (a) et (b) fait partie de la
définition \ref{definition-etale-smooth-local-source-target}.
L'équivalence de (a) et (e) est donnée par la
partie (2) du lemme \ref{lemma-etale-smooth-local-source-target-implies}.
L'équivalence de (a) et (f) est donnée par la
partie (1) du lemme \ref{lemma-etale-smooth-local-source-target-implies}.
Puisque (a) est maintenant équivalente à (e) et à (f), il s'ensuit que
(a) est équivalente à (g).

\medskip\noindent
Il est clair que (c) implique (b). Si (b) est vérifiée, alors, pour tout
$x \in |X|$, nous pouvons choisir un morphisme lisse
$b_x : V_x \to Y$, un morphisme \'etale $U_x \to V_x \times_Y X$,
et un point $u_x \in |U_x|$ qui s'envoie sur $x$,
de telle sorte que $U_x \to V_x$ vérifie $\mathcal{P}$.
Alors $h = \coprod h_x : \coprod U_x \to \coprod V_x$
forme, avec $a = \coprod a_x$ et $b = \coprod b_x$, un diagramme comme en (c).
(Remarquons que $h$ vérifie la propriété $\mathcal{P}$, puisque $\{V_x \to \coprod V_x\}$
est un recouvrement lisse et que $\mathcal{P}$ est locale pour la topologie lisse sur le but.)
Ainsi, (b) et (c) sont équivalentes.

\medskip\noindent
Nous savons maintenant que (a), (b), (c), (e), (f) et (g) sont équivalentes.
Supposons (a) vérifiée. Soient $U, V, a, b, h$ comme en (d). Alors
$X \times_Y V \to V$ vérifie $\mathcal{P}$ puisque $\mathcal{P}$ est stable par
changement de base lisse ; il s'ensuit que $U \to V$ vérifie $\mathcal{P}$ puisque $\mathcal{P}$
est stable par précomposition par des morphismes \'etales. Réciproquement, si (d)
est vérifiée, prendre $U = X$ et $V = Y$ montre que $f$ vérifie $\mathcal{P}$.
\end{proof}

\begin{lemma}
\label{lemma-etale-smooth-local-source-target}
Soit $S$ un schéma.
Soit $\mathcal{P}$ une propriété des morphismes d'espaces algébriques sur $S$.
Supposons que
\begin{enumerate}
\item $\mathcal{P}$ soit locale pour la topologie \'etale sur la source,
\item $\mathcal{P}$ soit locale pour la topologie lisse sur le but, et
\item $\mathcal{P}$ soit stable par postcomposition par des immersions ouvertes :
si $f : X \to Y$ vérifie $\mathcal{P}$ et si $Y \subset Z$ est une immersion ouverte,
alors $X \to Z$ vérifie $\mathcal{P}$.
\end{enumerate}
Alors $\mathcal{P}$ est locale pour la topologie \'etale sur la source et pour la topologie lisse sur le but.
\end{lemma}

\begin{proof}
Soit $\mathcal{P}$ une propriété des morphismes d'espaces algébriques qui
satisfait aux conditions (1), (2) et (3) du lemme. D'après le
lemme \ref{lemma-precompose-property-local-source},
on voit que $\mathcal{P}$ est stable par précomposition par des
morphismes \'etales. D'après le
lemme \ref{lemma-pullback-property-local-target},
on voit que $\mathcal{P}$ est stable par changement de base lisse.
Il suffit donc de prouver que la partie (3) de la
définition \ref{definition-local-source-target}
est vérifiée.

\medskip\noindent
Plus précisément, supposons que $f : X \to Y$ soit un morphisme
d'espaces algébriques sur $S$ qui vérifie la
partie (3)(b) de la définition \ref{definition-local-source-target}.
Autrement dit, pour tout $x \in X$, il existe
un morphisme lisse $b_x : V_x \to Y$,
un morphisme \'etale $U_x \to V_x \times_Y X$ et
un point $u_x \in |U_x|$ qui s'envoie sur $x$
de telle sorte que $h_x : U_x \to V_x$ vérifie $\mathcal{P}$.
Pour achever la preuve du lemme, il suffit de montrer que $f$ vérifie $\mathcal{P}$.

\medskip\noindent
Soit $a_x : U_x \to X$ le composé $U_x \to V_x \times_Y X \to X$.
Posons $U = \coprod U_x$, $a = \coprod a_x$, $V = \coprod V_x$,
$b = \coprod b_x$ et $h = \coprod h_x$. Nous obtenons un
diagramme commutatif
$$
\xymatrix{
U \ar[d]_a \ar[r]_h & V \ar[d]^b \\
X \ar[r]^f & Y
}
$$
où $b$ est lisse, $U \to V \times_Y X$ est \'etale et $a$ est surjectif.
Remarquons que $h$ vérifie $\mathcal{P}$ puisque chaque $h_x$ la vérifie et que $\mathcal{P}$
est locale pour la topologie lisse sur le but. Au paragraphe suivant, nous montrons
que l'on peut supposer que $U, V, X, Y$ sont des schémas ; nous invitons le lecteur
à passer ce paragraphe.

\medskip\noindent
Soient $X, Y, U, V, a, b, f, h$ comme au paragraphe précédent. Il faut
montrer que $f$ vérifie $\mathcal{P}$. Soit $X' \to X$ un morphisme \'etale surjectif
avec $X_i$ un schéma. Posons $U' = X' \times_X U$. Alors
$U' \to X'$ est surjectif et $U' \to X' \times_Y V$ est \'etale.
Comme $\mathcal{P}$ est locale pour la topologie \'etale sur la source, on voit que
$U' \to V$ vérifie $\mathcal{P}$ et qu'il suffit de montrer que
$X' \to Y$ vérifie $\mathcal{P}$. Autrement dit, on peut supposer
que $X$ est un schéma. Choisissons ensuite un morphisme \'etale surjectif
$Y' \to Y$, où $Y'$ est un schéma. Posons $V' = V \times_Y Y'$,
$X' = X \times_Y Y'$ et $U' = U \times_Y Y'$. Alors
$U' \to X'$ est surjectif et $U' \to X' \times_{Y'} V'$ est \'etale.
Comme $\mathcal{P}$ est locale pour la topologie lisse sur le but, on voit que $U' \to V'$ vérifie
$\mathcal{P}$ et qu'il suffit de prouver que $X' \to Y'$ vérifie $\mathcal{P}$.
Ainsi, on peut supposer que $X$ et $Y$ sont tous deux des schémas.
Choisissons un morphisme \'etale surjectif $V' \to V$
avec $V'$ un schéma. Posons $U' = U \times_V V'$.
Alors $U' \to X$ est surjectif et $U' \to X \times_Y V'$ est \'etale.
Comme $\mathcal{P}$ est locale pour la topologie lisse sur la source, on voit que
$U' \to V'$ vérifie $\mathcal{P}$. Nous pouvons donc remplacer $U, V$ par
$U', V'$ et supposer que $X, Y, V$ sont des schémas.
Enfin, remplaçons $U$ par un schéma muni d'un morphisme \'etale surjectif vers $U$ ;
on voit alors que l'on peut supposer que $U, V, X, Y$ sont tous des schémas.

\medskip\noindent
Si $U, V, X, Y$ sont des schémas, alors $f$ vérifie $\mathcal{P}$
d'après Descente, lemme \ref{descent-lemma-etale-tau-local-source-target}.
\end{proof}

\begin{remark}
\label{remark-list-etale-smooth-local-source-target}
À l'aide du lemme \ref{lemma-etale-smooth-local-source-target}
et des résultats établis dans les sections précédentes de ce chapitre, on dresse aisément
une liste de types de morphismes dont la propriété est locale pour la topologie lisse sur la
source et le but. Dans chaque cas, nous indiquons le lemme qui implique que
la propriété est locale pour la topologie étale sur la source et celui qui implique qu'elle
est locale pour la topologie lisse sur le but. Dans chaque cas, la troisième hypothèse
du lemme \ref{lemma-etale-smooth-local-source-target}
se vérifie immédiatement ; nous l'omettons. Voici la liste :
\begin{enumerate}
\item \'etale, voir les
lemmes \ref{lemma-etale-etale-local-source} et
\ref{lemma-descending-property-etale},
\item localement quasi-fini, voir les
lemmes \ref{lemma-locally-quasi-finite-etale-local-source} et
\ref{lemma-descending-property-quasi-finite},
\item non ramifié, voir les
lemmes \ref{lemma-unramified-etale-local-source} et
\ref{lemma-descending-property-unramified}, et
\item ajouter ici d'autres cas au besoin.
\end{enumerate}
Bien entendu, toute propriété énumérée dans la
remarque \ref{remark-list-local-source-target}
est a fortiori un exemple que l'on pourrait faire figurer ici.
\end{remark}






% CH074 slice boundary 3208
\section{Données de descente pour des espaces au-dessus d'espaces}
\label{section-descent-datum}

\noindent
Cette section est l'analogue de Descente, section
\ref{descent-section-descent-datum}, dans le cas des espaces algébriques.
La plupart des arguments de cette section sont formels et ne reposent
que sur la définition d'une donnée de descente.

\begin{definition}
\label{definition-descent-datum}
Soit $S$ un schéma. Soit $f : Y \to X$ un morphisme d'espaces algébriques
au-dessus de $S$.
\begin{enumerate}
\item Soit $V \to Y$ un morphisme d'espaces algébriques.
Une {\it donnée de descente pour $V/Y/X$} est un isomorphisme
$\varphi : V \times_X Y \to Y \times_X V$ d'espaces algébriques au-dessus de
$Y \times_X Y$ satisfaisant à la {\it condition de cocycle} suivante : le diagramme
$$
\xymatrix{
V \times_X Y \times_X Y \ar[rd]^{\varphi_{01}} \ar[rr]_{\varphi_{02}} &
&
Y \times_X Y \times_X V\\
&
Y \times_X V \times_X Y \ar[ru]^{\varphi_{12}}
}
$$
est commutatif (avec les notations évidentes).
\item On dit également que le couple $(V/Y, \varphi)$ est
une {\it donnée de descente relativement à $Y \to X$}.
\item Un {\it morphisme $f : (V/Y, \varphi) \to (V'/Y, \varphi')$ de données
de descente relativement à $Y \to X$} est un morphisme
$f : V \to V'$ d'espaces algébriques au-dessus de $Y$ tel que
le diagramme
$$
\xymatrix{
V \times_X Y \ar[r]_{\varphi} \ar[d]_{f \times \text{id}_Y} &
Y \times_X V \ar[d]^{\text{id}_Y \times f} \\
V' \times_X Y \ar[r]^{\varphi'} & Y \times_X V'
}
$$
soit commutatif.
\end{enumerate}
\end{definition}

\begin{remark}
\label{remark-easier}
Soit $S$ un schéma.
Soit $Y \to X$ un morphisme d'espaces algébriques au-dessus de $S$.
Soit $(V/Y, \varphi)$ une donnée de descente relativement à $Y \to X$.
On peut considérer l'isomorphisme $\varphi$ comme un isomorphisme
$$
(Y \times_X Y) \times_{\text{pr}_0, Y} V
\longrightarrow
(Y \times_X Y) \times_{\text{pr}_1, Y} V
$$
d'espaces algébriques au-dessus de $Y \times_X Y$. Ainsi, de manière informelle, on peut
considérer $\varphi$ comme une application
$\varphi : \text{pr}_0^*V \to \text{pr}_1^*V$\footnote{Malheureusement,
nous avons choisi ici le « mauvais » sens pour notre flèche. Dans les
définitions \ref{definition-descent-datum} et
\ref{definition-descent-datum-for-family-of-morphisms}
il faudrait prendre le sens opposé à celui adopté dans la
définition \ref{definition-descent-datum-quasi-coherent},
selon le principe général que les « fonctions » et les « espaces » sont duaux.}.
La condition de cocycle exprime alors que
$\text{pr}_{02}^*\varphi =
\text{pr}_{12}^*\varphi \circ \text{pr}_{01}^*\varphi$.
Elle est ainsi très semblable à celle d'une donnée de descente sur des
faisceaux quasi-cohérents.
\end{remark}

\noindent
Voici la définition lorsqu'on dispose d'une famille de morphismes
à cible fixée.

\begin{definition}
\label{definition-descent-datum-for-family-of-morphisms}
Soit $S$ un schéma.
Soit $\{X_i \to X\}_{i \in I}$ une famille de morphismes
 d'espaces algébriques au-dessus de $S$, à cible fixée $X$.
\begin{enumerate}
\item Une {\it donnée de descente $(V_i, \varphi_{ij})$ relativement à la
famille $\{X_i \to X\}$} est constituée d'un espace algébrique $V_i$ au-dessus de $X_i$
pour chaque $i \in I$, et d'un isomorphisme
$\varphi_{ij} : V_i \times_X X_j \to X_i \times_X V_j$
d'espaces algébriques au-dessus de $X_i \times_X X_j$, pour chaque couple $(i, j) \in I^2$,
tel que, pour tout triplet d'indices $(i, j, k) \in I^3$,
le diagramme
$$
\xymatrix{
V_i \times_X X_j \times_X X_k
\ar[rd]^{\text{pr}_{01}^*\varphi_{ij}}
\ar[rr]_{\text{pr}_{02}^*\varphi_{ik}} &
&
X_i \times_X X_j \times_X V_k\\
&
X_i \times_X V_j \times_X X_k
\ar[ru]^{\text{pr}_{12}^*\varphi_{jk}}
}
$$
d'espaces algébriques au-dessus de $X_i \times_X X_j \times_X X_k$ soit commutatif
(avec les notations évidentes).
\item Un {\it morphisme
$\psi : (V_i, \varphi_{ij}) \to (V'_i, \varphi'_{ij})$
de données de descente} est constitué d'une famille $\psi = (\psi_i)_{i \in I}$
de morphismes $\psi_i : V_i \to V'_i$ d'espaces algébriques au-dessus de $X_i$
telle que tous les diagrammes
$$
\xymatrix{
V_i \times_X X_j \ar[r]_{\varphi_{ij}} \ar[d]_{\psi_i \times \text{id}} &
X_i \times_X V_j \ar[d]^{\text{id} \times \psi_j} \\
V'_i \times_X X_j \ar[r]^{\varphi'_{ij}} & X_i \times_X V'_j
}
$$
soient commutatifs.
\end{enumerate}
\end{definition}

\begin{remark}
\label{remark-easier-family}
Soit $S$ un schéma.
Soit $\{X_i \to X\}_{i \in I}$ une famille de morphismes
 d'espaces algébriques au-dessus de $S$, à cible fixée $X$.
Soit $(V_i, \varphi_{ij})$ une donnée de descente relativement à
$\{X_i \to X\}$. On peut considérer les isomorphismes $\varphi_{ij}$
comme des isomorphismes
$$
(X_i \times_X X_j) \times_{\text{pr}_0, X_i} V_i
\longrightarrow
(X_i \times_X X_j) \times_{\text{pr}_1, X_j} V_j
$$
d'espaces algébriques au-dessus de $X_i \times_X X_j$. Ainsi, de manière informelle, on peut
considérer $\varphi_{ij}$ comme un isomorphisme
$\text{pr}_0^*V_i \to \text{pr}_1^*V_j$ au-dessus de $X_i \times_X X_j$.
La condition de cocycle exprime alors que
$\text{pr}_{02}^*\varphi_{ik} =
\text{pr}_{12}^*\varphi_{jk} \circ \text{pr}_{01}^*\varphi_{ij}$.
Elle est ainsi très semblable à celle d'une donnée de descente sur des
faisceaux quasi-cohérents.
\end{remark}

\noindent
La raison pour laquelle nous travaillerons habituellement avec une famille réduite
à un seul morphisme est le lemme suivant.

\begin{lemma}
\label{lemma-family-is-one}
Soit $S$ un schéma.
Soit $\{X_i \to X\}_{i \in I}$ une famille de morphismes
 d'espaces algébriques au-dessus de $S$, à cible fixée $X$.
Posons $Y = \coprod_{i \in I} X_i$.
Il existe une équivalence canonique de catégories
$$
\begin{matrix}
\text{catégorie des données de descente } \\
\text{relativement à la famille } \{X_i \to X\}_{i \in I}
\end{matrix}
\longrightarrow
\begin{matrix}
\text{ catégorie des données de descente} \\
\text{ relativement à } Y/X
\end{matrix}
$$
qui envoie $(V_i, \varphi_{ij})$ sur $(V, \varphi)$, où
$V = \coprod_{i\in I} V_i$ et $\varphi = \coprod \varphi_{ij}$.
\end{lemma}

\begin{proof}
Remarquons que $Y \times_X Y = \coprod_{ij} X_i \times_X X_j$
et de même pour les produits fibrés itérés.
Se donner un morphisme $V \to Y$ revient exactement à
se donner une famille $V_i \to X_i$. De même, se donner une donnée de descente
$\varphi$ revient exactement à se donner une famille $\varphi_{ij}$.
\end{proof}

\begin{lemma}
\label{lemma-pullback}
Changement de base des données de descente. Soit $S$ un schéma.
\begin{enumerate}
\item Soit
$$
\xymatrix{
Y' \ar[r]_f \ar[d]_{a'} & Y \ar[d]^a \\
X' \ar[r]^h & X
}
$$
soit un diagramme commutatif d'espaces algébriques au-dessus de $S$.
La construction
$$
(V \to Y, \varphi) \longmapsto f^*(V \to Y, \varphi) = (V' \to Y', \varphi')
$$
où $V' = Y' \times_Y V$ et où
$\varphi'$ est définie comme la composée
$$
\xymatrix{
V' \times_{X'} Y' \ar@{=}[r] &
(Y' \times_Y V) \times_{X'} Y' \ar@{=}[r] &
(Y' \times_{X'} Y') \times_{Y \times_X Y} (V \times_X Y)
\ar[d]^{\text{id} \times \varphi} \\
Y' \times_{X'} V' \ar@{=}[r] &
Y' \times_{X'} (Y' \times_Y V) &
(Y' \times_X Y') \times_{Y \times_X Y} (Y \times_X V) \ar@{=}[l]
}
$$
définit un foncteur de la catégorie des données de descente
relativement à $Y \to X$ vers la catégorie des données de descente
relativement à $Y' \to X'$.
\item Étant donnés deux morphismes $f_i : Y' \to Y$, $i = 0, 1$, rendant le
diagramme commutatif, les foncteurs $f_0^*$ et $f_1^*$ sont
canoniquement isomorphes.
\end{enumerate}
\end{lemma}

\begin{proof}
Nous omettons la démonstration de (1), mais remarquons que le morphisme
$\varphi'$ est le morphisme $(f \times f)^*\varphi$ dans les
notations introduites dans la remarque \ref{remark-easier}.
Pour (2), indiquons quel morphisme
$f_0^*V \to f_1^*V$ fournit l'isomorphisme fonctoriel. En effet,
puisque $f_0$ et $f_1$ s'insèrent tous deux dans le diagramme commutatif,
il existe un unique morphisme $r : Y' \to Y \times_X Y$
tel que $f_i = \text{pr}_i \circ r$. Prenons alors
\begin{eqnarray*}
f_0^*V & = &
Y' \times_{f_0, Y} V \\
& = &
Y' \times_{\text{pr}_0 \circ r, Y} V \\
& = &
Y' \times_{r, Y \times_X Y} (Y \times_X Y) \times_{\text{pr}_0, Y} V \\
& \xrightarrow{\varphi} &
Y' \times_{r, Y \times_X Y} (Y \times_X Y) \times_{\text{pr}_1, Y} V \\
& = &
Y' \times_{\text{pr}_1 \circ r, Y} V \\
& = &
Y' \times_{f_1, Y} V \\
& = & f_1^*V
\end{eqnarray*}
Nous omettons la vérification.
\end{proof}

\begin{definition}
\label{definition-pullback-functor}
Avec $S, X, X', Y, Y', f, a, a', h$ comme dans le lemme \ref{lemma-pullback},
le foncteur
$$
(V, \varphi) \longmapsto f^*(V, \varphi)
$$
construit dans ce lemme est appelé le {\it foncteur de changement de base} des données de descente.
\end{definition}

\begin{lemma}
\label{lemma-pullback-family}
Soit $S$ un schéma. Soient $\mathcal{U}' = \{X'_i \to X'\}_{i \in I'}$ et
$\mathcal{U} = \{X_j \to X\}_{i \in I}$ des familles de morphismes à
cible fixée. Soient $\alpha : I' \to I$, $g : X' \to X$ et les morphismes
$g_i : X'_i \to X_{\alpha(i)}$ les composantes d'un morphisme de familles
de morphismes à cible fixée, voir
Sites, définition \ref{sites-definition-morphism-coverings}.
\begin{enumerate}
\item Soit $(V_i, \varphi_{ij})$ une donnée de descente relativement à la
famille $\mathcal{U}$. Le système
$$
\left(
g_i^*V_{\alpha(i)}, (g_i \times g_j)^*\varphi_{\alpha(i) \alpha(j)}
\right)
$$
(avec les notations de la remarque \ref{remark-easier-family})
est une donnée de descente relativement à $\mathcal{U}'$.
\item Cette construction définit un foncteur de la catégorie des
données de descente relativement à $\mathcal{U}$ vers la catégorie des
données de descente relativement à $\mathcal{U}'$.
\item Soient aussi $\beta : I' \to I$, $h : X' \to X$ et les morphismes
$h'_i : X'_i \to X_{\beta(i)}$ les composantes d'un second morphisme de familles
de morphismes à cible fixée. Si $g = h$, alors les deux foncteurs obtenus
entre les catégories de données de descente sont canoniquement isomorphes.
\item Via le lemme \ref{lemma-family-is-one}, ces foncteurs coïncident
avec les foncteurs de changement de base construits dans le lemme \ref{lemma-pullback}.
\end{enumerate}
\end{lemma}

\begin{proof}
L'assertion résulte du lemme \ref{lemma-pullback} au moyen de la
correspondance du lemme \ref{lemma-family-is-one}.
\end{proof}

\begin{definition}
\label{definition-pullback-functor-family}
Avec $\mathcal{U}' = \{X'_i \to X'\}_{i \in I'}$,
$\mathcal{U} = \{X_i \to X\}_{i \in I}$, $\alpha : I' \to I$,
$g : X' \to X$ et $g_i : X'_i \to X_{\alpha(i)}$ comme dans le
lemme \ref{lemma-pullback-family}, le foncteur
$$
(V_i, \varphi_{ij}) \longmapsto
(g_i^*V_{\alpha(i)}, (g_i \times g_j)^*\varphi_{\alpha(i) \alpha(j)})
$$
construit dans ce lemme
est appelé le {\it foncteur de changement de base} des données de descente.
\end{definition}

\noindent
Si $\mathcal{U}$ et $\mathcal{U}'$ ont la même cible $X$,
et si $\mathcal{U}'$ raffine $\mathcal{U}$ (voir
Sites, définition \ref{sites-definition-morphism-coverings}),
mais qu'aucun couple explicite $(\alpha, g_i)$ n'est donné, nous pouvons néanmoins
parler du foncteur de changement de base, puisque le
lemme \ref{lemma-pullback-family} montre que le choix du couple n'importe pas
(à isomorphisme canonique près).

\begin{definition}
\label{definition-effective}
Soit $S$ un schéma. Soit $f : Y \to X$ un morphisme d'espaces algébriques au-dessus de
$S$.
\begin{enumerate}
\item À tout espace algébrique $U$ au-dessus de $X$ est associée la
{\it donnée de descente triviale} de $U$ relativement à $\text{id} : X \to X$, à savoir
le morphisme identité de $U$.
\item Le lemme \ref{lemma-pullback} fournit alors une
{\it donnée de descente canonique} sur $Y \times_X U$
relativement à $Y \to X$, obtenue par changement de base de la donnée de descente
triviale par $f$. Cette donnée de descente est souvent
notée $(Y \times_X U, can)$.
\item Une donnée de descente $(V, \varphi)$ relativement à $Y/X$
est dite {\it effective} si $(V, \varphi)$
est isomorphe à la donnée de descente canonique
$(Y \times_X U, can)$ pour un espace algébrique $U$ au-dessus de $X$.
\end{enumerate}
\end{definition}

\noindent
Ainsi, l'effectivité signifie qu'il existe un espace algébrique $U$
au-dessus de $X$ et un isomorphisme $\psi : V \to Y \times_X U$
au-dessus de $Y$ tels que $\varphi$ soit égale à la composée
$$
V \times_X Y \xrightarrow{\psi \times \text{id}_Y}
Y \times_X U \times_S Y =
Y \times_X Y \times_X U
\xrightarrow{\text{id}_Y \times \psi^{-1}}
Y \times_X V
$$
Il existe ici une légère difficulté : cette définition
(dans son esprit) entre en conflit avec celle donnée dans
Descente, définition \ref{descent-definition-effective},
lorsque $Y$ et $X$ sont des schémas. Cependant, le contexte permettra toujours
de savoir laquelle des deux est visée.

\begin{definition}
\label{definition-effective-family}
Soit $S$ un schéma.
Soit $\{X_i \to X\}$ une famille de morphismes d'espaces algébriques au-dessus de $S$,
à cible fixée $X$.
\begin{enumerate}
\item Étant donné un espace algébrique $U$ au-dessus de $X$,
il existe une {\it donnée de descente canonique} sur la famille des
espaces algébriques $X_i \times_X U$, obtenue par changement de base de la donnée de descente
triviale de $U$ relativement à $\{\text{id} : S \to S\}$.
Nous notons cette donnée de descente $(X_i \times_X U, can)$.
\item Une donnée de descente $(V_i, \varphi_{ij})$
relativement à $\{X_i \to S\}$ est dite {\it effective}
s'il existe un espace algébrique $U$ au-dessus de $X$ tel que
$(V_i, \varphi_{ij})$ soit isomorphe à $(X_i \times_X U, can)$.
\end{enumerate}
\end{definition}






\section{Données de descente en termes de faisceaux}
\label{section-descent-data-sheaves}

\noindent
Cette section est l'analogue de Descente, section
\ref{descent-section-descent-data-sheaves}.
Elle diffère légèrement, puisque les espaces algébriques sont déjà des faisceaux.

\begin{lemma}
\label{lemma-descent-data-sheaves}
Soit $S$ un schéma. Soit $\{X_i \to X\}_{i \in I}$ un recouvrement fppf
d'espaces algébriques au-dessus de $S$ (Topologies sur les espaces,
définition \ref{spaces-topologies-definition-fppf-covering}).
Il existe une équivalence de catégories
$$
\left\{
\begin{matrix}
\text{données de descente }(V_i, \varphi_{ij})\\
\text{relativement à }\{X_i \to X\}
\end{matrix}
\right\}
\leftrightarrow
\left\{
\begin{matrix}
\text{faisceaux }F\text{ sur }(\Sch/S)_{fppf}\text{ munis}\\
\text{d'un morphisme }F \to X\text{ tel que chaque}\\
X_i \times_X F\text{ soit un espace algébrique}
\end{matrix}
\right\}.
$$
De plus,
\begin{enumerate}
\item l'espace algébrique $X_i \times_X F$ du membre de droite
correspond à $V_i$ du membre de gauche, et
\item le faisceau $F$ est un espace algébrique\footnote{Nous verrons
plus loin que c'est toujours le cas si $I$ n'est pas trop grand, voir
Amorçage, lemme \ref{bootstrap-lemma-descend-algebraic-space}.}
si et seulement si la
donnée de descente $(X_i, \varphi_{ij})$ correspondante est effective.
\end{enumerate}
\end{lemma}

\begin{proof}
Construisons le foncteur de droite à gauche.
Soit $F \to X$ un morphisme de faisceaux sur $(\Sch/S)_{fppf}$ tel que chaque
$V_i = X_i \times_X F$ soit un espace algébrique. Nous disposons de la
projection $V_i \to X_i$.
Alors $V_i \times_X X_j$ et $X_i \times_X V_j$
représentent tous deux le faisceau $X_i \times_X F \times_X X_j$ ;
nous obtenons donc un isomorphisme
$$
\varphi_{ii'} : V_i \times_X X_j \to X_i \times_X V_j
$$
Il est immédiat que les applications $\varphi_{ij}$
sont des morphismes au-dessus de $X_i \times_X X_j$ et satisfont à la
condition de cocycle. Le foncteur de droite à gauche est donné
par cette construction $F \mapsto (V_i, \varphi_{ij})$.

\medskip\noindent
Construisons maintenant un foncteur de gauche à droite.
Les isomorphismes $\varphi_{ij}$ fournissent des isomorphismes
$$
\varphi_{ij} : V_i \times_X X_j \longrightarrow X_i \times_X V_j
$$
au-dessus de $X_i \times X_j$. Posons $F$ égal au coégalisateur du
diagramme suivant :
$$
\xymatrix{
\coprod_{i, i'} V_i \times_X X_j
\ar@<1ex>[rr]^-{\text{pr}_0}
\ar@<-1ex>[rr]_-{\text{pr}_1 \circ \varphi_{ij}}
& &
\coprod_i V_i \ar[r]
&
F
}
$$
La condition de cocycle assure que $F$ est muni d'un morphisme
$F \to X$ et que $X_i \times_X F$ est isomorphe à $V_i$.
Le foncteur de gauche à droite est donné
par cette construction $(V_i, \varphi_{ij}) \mapsto F$.

\medskip\noindent
Nous omettons de vérifier que ces constructions
sont des foncteurs quasi-inverses l'un de l'autre. Les derniers énoncés
(1) et (2) résultent des constructions.
\end{proof}




\begin{multicols}{2}[\section{Autres chapitres}]
\noindent
Préliminaires
\begin{enumerate}
\item \hyperref[introduction-section-phantom]{Introduction}
\item \hyperref[conventions-section-phantom]{Conventions}
\item \hyperref[sets-section-phantom]{Théorie des ensembles}
\item \hyperref[categories-section-phantom]{Catégories}
\item \hyperref[topology-section-phantom]{Topologie}
\item \hyperref[sheaves-section-phantom]{Faisceaux sur les espaces}
\item \hyperref[sites-section-phantom]{Sites et faisceaux}
\item \hyperref[stacks-section-phantom]{Champs}
\item \hyperref[fields-section-phantom]{Corps}
\item \hyperref[algebra-section-phantom]{Algèbre commutative}
\item \hyperref[brauer-section-phantom]{Groupes de Brauer}
\item \hyperref[homology-section-phantom]{Algèbre homologique}
\item \hyperref[derived-section-phantom]{Catégories dérivées}
\item \hyperref[simplicial-section-phantom]{Méthodes simpliciales}
\item \hyperref[more-algebra-section-phantom]{Compléments d'algèbre}
\item \hyperref[smoothing-section-phantom]{Lissage des morphismes d'anneaux}
\item \hyperref[modules-section-phantom]{Faisceaux de modules}
\item \hyperref[sites-modules-section-phantom]{Modules sur les sites}
\item \hyperref[injectives-section-phantom]{Objets injectifs}
\item \hyperref[cohomology-section-phantom]{Cohomologie des faisceaux}
\item \hyperref[sites-cohomology-section-phantom]{Cohomologie sur les sites}
\item \hyperref[dga-section-phantom]{Algèbre différentielle graduée}
\item \hyperref[dpa-section-phantom]{Algèbre à puissances divisées}
\item \hyperref[sdga-section-phantom]{Faisceaux différentiels gradués}
\item \hyperref[hypercovering-section-phantom]{Hyperrecouvrements}
\end{enumerate}
Schémas
\begin{enumerate}
\setcounter{enumi}{25}
\item \hyperref[schemes-section-phantom]{Schémas}
\item \hyperref[constructions-section-phantom]{Constructions de schémas}
\item \hyperref[properties-section-phantom]{Propriétés des schémas}
\item \hyperref[morphisms-section-phantom]{Morphismes de schémas}
\item \hyperref[coherent-section-phantom]{Cohomologie des schémas}
\item \hyperref[divisors-section-phantom]{Diviseurs}
\item \hyperref[limits-section-phantom]{Limites de schémas}
\item \hyperref[varieties-section-phantom]{Variétés}
\item \hyperref[topologies-section-phantom]{Topologies sur les schémas}
\item \hyperref[descent-section-phantom]{Descente}
\item \hyperref[perfect-section-phantom]{Catégories dérivées de schémas}
\item \hyperref[more-morphisms-section-phantom]{Compléments sur les morphismes}
\item \hyperref[flat-section-phantom]{Compléments sur la platitude}
\item \hyperref[groupoids-section-phantom]{Schémas en groupoïdes}
\item \hyperref[more-groupoids-section-phantom]{Compléments sur les schémas en groupoïdes}
\item \hyperref[etale-section-phantom]{Morphismes étales de schémas}
\end{enumerate}
Thèmes de théorie des schémas
\begin{enumerate}
\setcounter{enumi}{41}
\item \hyperref[chow-section-phantom]{Homologie de Chow}
\item \hyperref[intersection-section-phantom]{Théorie de l'intersection}
\item \hyperref[pic-section-phantom]{Schémas de Picard des courbes}
\item \hyperref[weil-section-phantom]{Théories cohomologiques de Weil}
\item \hyperref[adequate-section-phantom]{Modules adéquats}
\item \hyperref[dualizing-section-phantom]{Complexes dualisants}
\item \hyperref[duality-section-phantom]{Dualité pour les schémas}
\item \hyperref[discriminant-section-phantom]{Discriminants et différentes}
\item \hyperref[derham-section-phantom]{Cohomologie de de Rham}
\item \hyperref[local-cohomology-section-phantom]{Cohomologie locale}
\item \hyperref[algebraization-section-phantom]{Géométrie algébrique et formelle}
\item \hyperref[curves-section-phantom]{Courbes algébriques}
\item \hyperref[resolve-section-phantom]{Résolution des surfaces}
\item \hyperref[models-section-phantom]{Réduction semi-stable}
\item \hyperref[functors-section-phantom]{Foncteurs et morphismes}
\item \hyperref[equiv-section-phantom]{Catégories dérivées de variétés}
\item \hyperref[pione-section-phantom]{Groupes fondamentaux de schémas}
\item \hyperref[etale-cohomology-section-phantom]{Cohomologie étale}
\item \hyperref[crystalline-section-phantom]{Cohomologie cristalline}
\item \hyperref[proetale-section-phantom]{Cohomologie pro-étale}
\item \hyperref[relative-cycles-section-phantom]{Cycles relatifs}
\item \hyperref[more-etale-section-phantom]{Compléments de cohomologie étale}
\item \hyperref[trace-section-phantom]{La formule des traces}
\end{enumerate}
Espaces algébriques
\begin{enumerate}
\setcounter{enumi}{64}
\item \hyperref[spaces-section-phantom]{Espaces algébriques}
\item \hyperref[spaces-properties-section-phantom]{Propriétés des espaces algébriques}
\item \hyperref[spaces-morphisms-section-phantom]{Morphismes d'espaces algébriques}
\item \hyperref[decent-spaces-section-phantom]{Espaces algébriques décents}
\item \hyperref[spaces-cohomology-section-phantom]{Cohomologie des espaces algébriques}
\item \hyperref[spaces-limits-section-phantom]{Limites d'espaces algébriques}
\item \hyperref[spaces-divisors-section-phantom]{Diviseurs sur les espaces algébriques}
\item \hyperref[spaces-over-fields-section-phantom]{Espaces algébriques sur des corps}
\item \hyperref[spaces-topologies-section-phantom]{Topologies sur les espaces algébriques}
\item \hyperref[spaces-descent-section-phantom]{Descente et espaces algébriques}
\item \hyperref[spaces-perfect-section-phantom]{Catégories dérivées d'espaces}
\item \hyperref[spaces-more-morphisms-section-phantom]{Compléments sur les morphismes d'espaces}
\item \hyperref[spaces-flat-section-phantom]{Platitude sur les espaces algébriques}
\item \hyperref[spaces-groupoids-section-phantom]{Groupoïdes dans les espaces algébriques}
\item \hyperref[spaces-more-groupoids-section-phantom]{Compléments sur les groupoïdes dans les espaces}
\item \hyperref[bootstrap-section-phantom]{Amorçage}
\item \hyperref[spaces-pushouts-section-phantom]{Sommes amalgamées d'espaces algébriques}
\end{enumerate}
Thèmes de géométrie
\begin{enumerate}
\setcounter{enumi}{81}
\item \hyperref[spaces-chow-section-phantom]{Groupes de Chow des espaces}
\item \hyperref[groupoids-quotients-section-phantom]{Quotients de groupoïdes}
\item \hyperref[spaces-more-cohomology-section-phantom]{Compléments sur la cohomologie des espaces}
\item \hyperref[spaces-simplicial-section-phantom]{Espaces simpliciaux}
\item \hyperref[spaces-duality-section-phantom]{Dualité pour les espaces}
\item \hyperref[formal-spaces-section-phantom]{Espaces algébriques formels}
\item \hyperref[restricted-section-phantom]{Algébrisation des espaces formels}
\item \hyperref[spaces-resolve-section-phantom]{Résolution des surfaces revisitée}
\end{enumerate}
Théorie des déformations
\begin{enumerate}
\setcounter{enumi}{89}
\item \hyperref[formal-defos-section-phantom]{Théorie formelle des déformations}
\item \hyperref[defos-section-phantom]{Théorie des déformations}
\item \hyperref[cotangent-section-phantom]{Le complexe cotangent}
\item \hyperref[examples-defos-section-phantom]{Problèmes de déformation}
\end{enumerate}
Champs algébriques
\begin{enumerate}
\setcounter{enumi}{93}
\item \hyperref[algebraic-section-phantom]{Champs algébriques}
\item \hyperref[examples-stacks-section-phantom]{Exemples de champs}
\item \hyperref[stacks-sheaves-section-phantom]{Faisceaux sur les champs algébriques}
\item \hyperref[criteria-section-phantom]{Critères de représentabilité}
\item \hyperref[artin-section-phantom]{Axiomes d'Artin}
\item \hyperref[quot-section-phantom]{Espaces de Quot et de Hilbert}
\item \hyperref[stacks-properties-section-phantom]{Propriétés des champs algébriques}
\item \hyperref[stacks-morphisms-section-phantom]{Morphismes de champs algébriques}
\item \hyperref[stacks-limits-section-phantom]{Limites de champs algébriques}
\item \hyperref[stacks-cohomology-section-phantom]{Cohomologie des champs algébriques}
\item \hyperref[stacks-perfect-section-phantom]{Catégories dérivées de champs}
\item \hyperref[stacks-introduction-section-phantom]{Introduction aux champs algébriques}
\item \hyperref[stacks-more-morphisms-section-phantom]{Compléments sur les morphismes de champs}
\item \hyperref[stacks-geometry-section-phantom]{La géométrie des champs}
\end{enumerate}
Thèmes de théorie des espaces de modules
\begin{enumerate}
\setcounter{enumi}{107}
\item \hyperref[moduli-section-phantom]{Champs de modules}
\item \hyperref[moduli-curves-section-phantom]{Espaces de modules de courbes}
\end{enumerate}
Divers
\begin{enumerate}
\setcounter{enumi}{109}
\item \hyperref[examples-section-phantom]{Exemples}
\item \hyperref[exercises-section-phantom]{Exercices}
\item \hyperref[guide-section-phantom]{Guide bibliographique}
\item \hyperref[desirables-section-phantom]{Desiderata}
\item \hyperref[coding-section-phantom]{Style de programmation}
\item \hyperref[obsolete-section-phantom]{Obsolète}
\item \hyperref[fdl-section-phantom]{Licence de documentation libre GNU}
\item \hyperref[index-section-phantom]{Index généré automatiquement}
\end{enumerate}
\end{multicols}


\bibliography{my}
\bibliographystyle{amsalpha}


\end{document}
